\documentclass[../main.tex]{subfiles}
\begin{document}
%==========================================TIÊU ĐỀ TẬP========================================
\begin{table}[h]
\begin{adjustwidth}{-1cm}{}
    \begin{tabular}{>{\centering\arraybackslash}p{6cm}|>{\raggedright\arraybackslash}p{12.5cm}}
        \multirow{5}{6cm}
        % Cột bên trái: ÉPISODE và số tập
        {\\[-40pt]
            \flaregothic\fontsize{20pt}{20pt}\selectfont \centering
            \color{eptype}HISTOIRE\color{black}\\
            \freshbot\fontsize{72pt}{72pt}\selectfont
            \color{epnum}7\\[6pt]
            \cafeteria\fontsize{20pt}{20pt}\selectfont\color{epnum}(D-7)\color{black}
        \\[18pt]
        % \flaregothic\fontsize{14pt}{14pt}\selectfont
        % \color{epattr}BIENVENUE, \uline{\color{Banpire} Mel} !
        % \flaregothic\fontsize{18pt}{18pt}\selectfont
        % \color{epattr}ÉPISODE PERDU
        \color{black}
        }
        &
        % Dòng thứ nhất: tên tập tiếng Nhật
        {\fotskip\fontsize{18pt}{18pt}\selectfont \ruby{異}{い}\ruby{次}{じ}\ruby{元}{げん}の\ruby{血}{けつ}\ruby{脈}{みゃく}と\ruby{家}{か}\ruby{族}{ぞく}の\ruby{絆}{きずな}}
        \\[6pt]
        % Dòng thứ hai: tên tập tiếng Anh
        & {\cafeteria\fontsize{18pt}{18pt}\selectfont The Hikawa Invasion}\\[4pt]
        % Dòng thứ ba: tên tập tiếng Việt
        & {\vnmsans\fontsize{18pt}{18pt}\selectfont Cuộc đoàn tụ nhà Hikawa}\\[8pt]
        % òng thứ tư: Nhân vật xuất hiện
        & {\caviar{\fontsize{14pt}{14pt}\selectfont Nhân vật: \emp{kuon1} \emp{ryuuhou1} \emp{kaigou1} \emp{suzuno1} \emp{naomi1} \emp{game} \emp{coco} \emp{aloe}}}
        \\[8pt]
        % Dòng cuối cùng: Release Date
        % & {\continuum\fontsize{16pt}{16pt}\selectfont Release Date: 07/02/2021}\\[18pt]
        \end{tabular}
\end{adjustwidth}
\end{table}
%=========================================TRUYỆN CHÍNH========================================
\color{black}

\txtbx[2]
{{\color{vol} \uline{Tóm tắt:}} Trở về sau chuyến du hành xuyên qua ba thế giới song song đầy hiểm nguy, Kuon mang theo một ``món quà lưu niệm'' không ai ngờ tới: ba cô gái lạ mặt gồm Kaigou, Suzuno và Naomi. Cuộc gặp gỡ định mệnh tại nhà Hikawa tưởng chừng chỉ là sự đồng cảm giữa những linh hồn lạc lối, nhưng kết quả xét nghiệm ADN sáng hôm sau lại mở ra một sự thật chấn động. Dòng máu biết cách tìm về với nhau vượt qua mọi ranh giới không gian, biến người anh độc nhất vô nhị của dòng họ Hikawa thành anh cả của bốn cô em gái độc đáo.}

\begin{center}
    \uline{(Tiếp tục từ {\flaregothic HISTOIRE \dawnland 6}, quyển Khoảng Không.)}
\end{center}

\pov{Hikawa Kuon}

\lang{Etsunan}

Nếu có ai đó nói với tôi rằng một chuyến du hành thứ nguyên có thể biến một gã sinh viên lập trình cô đơn thành anh cả của một đại gia đình năm anh em, tôi sẽ cười vào mặt họ và bảo hãy bớt đọc light novel lại.

Nhưng cuộc sống vốn là kịch bản quái dị nhất mà không một nhà lập trình nào có thể viết trước mã nguồn, hay một nhà biên kịch phim nào có thể ngờ tới.

Chuyến hành trình băng qua ba thế giới song song - từ những tàn tích hỗn loạn của Nijisanji cho tới các biến cố tại vũ trụ \hll\ - cuối cùng cũng khép lại. Tôi bước ra khỏi khe nứt không gian với tư thế của một kẻ ``đi một về bốn'', hay đúng hơn là năm nếu tính cả Game - thực thể trí tuệ nhân tạo đang ẩn mình trong đồng hồ đeo tay của tôi. Phía sau lưng tôi là Kaigou, cô gái với mái tóc đuôi ngựa buộc lệch mang ánh mắt sắc lạnh như dao; Suzuno, cô bé có đôi mắt dị sắc mang nụ cười ngây thơ cùng chiếc vòng tay gắn ba quả chuông vàng; và Naomi, thiên thần nhỏ tám tuổi mang đôi mắt xanh lục bảo trong veo như nước hồ mùa thu.

Đúng như lời hứa, sau khi tôi tan ca lúc sáu giờ chiều, tôi dẫn ba cô gái về ký túc xá kiêm nhà Hikawa chúng tôi. Ý định ban đầu của tôi chỉ đơn thuần là giới thiệu họ như những ``món quà lưu niệm kỳ lạ'' từ chuyến du hành dài ngày, tìm cho họ một chỗ dừng chân an toàn sau biết bao giông bão. Trên đường đi, tôi đã dừng chân ở một cửa hàng tiện lợi mua đủ đồ ăn cho bữa tối, nhưng càng bước gần đến cửa nhà, tôi càng thấy lòng mình rối tung.

Vấn đề là, tôi nên giải thích thế nào đây? Làm sao mà nói với bố mẹ rằng đứa con trai độc nhất của họ đột nhiên mang về ba cô gái lạ, và muốn họ chấp nhận họ như một phần của gia đình?

Kaigou, sánh bước bên cạnh tôi, dường như đọc được sự bối rối trong mắt tôi, liền đá nhẹ vào chân tôi, giọng đầy châm chọc: ``Này, Kuon, bây giờ lại rụt rè nữa rồi? Hôm qua còn mạnh miệng nói sẽ giải thích mọi chuyện với bố mẹ mà.'' Tôi quay sang nhìn cô, chẹp miệng đáp: ``Cô thì biết cái đếch gì, từ trước đến nay có lúc nào tôi phải giải thích chuyện lớn thế này đâu.''

Suzuno, bước phía sau chúng tôi với chiếc vòng tay chuông vàng cứ leng keng theo mỗi bước chân, thấy hai người cãi nhau liền vội vàng kéo nhẹ áo tôi, giọng nhỏ nhẹ như gió thoảng: ``Anh chị đừng cãi nhau nữa ạ\dots\ Em cũng đã chuẩn bị tinh thần rồi, nếu bố mẹ không chấp nhận chúng em thì cũng không sao đâu ạ\dots'' 

Tôi vừa định an ủi em bé thì Naomi, lúc nào cũng lặng lẽ đứng bên cạnh Suzuno, đột nhiên nhỏ giọng kéo tay áo chị gái, mắt to tròn nhìn tôi nói: ``Anh à, Naomi ngửi thấy mùi thơm từ trong nhà ra rồi\dots\ Có mùi gì giống nấm không ạ?'' Câu nói của em nhỏ làm cả ba người chúng tôi đều dừng lại, đúng lúc cánh cửa gỗ quen thuộc bật mở ra.

Khi chúng tôi lên tới tầng năm - căn phòng quen thuộc của bố con tôi - bật mở, mùi thơm nức của nồi canh nấm nấu gừng từ bếp tỏa ra không gian. Mẹ tôi, tay vẫn còn vương mùi gừng tươi, bước ra từ nhà bếp. Bố tôi đang đứng bên kệ sách, chiếc áo len xám khoác hờ trên vai còn chưa kịp treo lên móc; dường như ông cũng vừa mới về tới nhà.

``Con về rồi đây bố mẹ ơi\dots'' tôi cất giọng chào quen thuộc.

Ngay khoảnh khắc ánh mắt chạm nhau, hai chị em tới từ vũ trụ song song bỗng đứng chết trân tại chỗ. Mắt họ mở to, đồng tử co lại như thể vừa bắt gặp một bức tranh thiêng liêng bước ra từ tiềm thức.

``\dots Chị Kaigou? Chị Suzuno?'' Naomi ngước nhìn hai cô chị đang khựng lại, ánh mắt mở to ngạc nhiên. ``Hai người sao vậy ạ?''

``Sao hai người lại đứng như trời trồng thế?'' tôi cũng không kìm được mà hỏi với sự tò mò nhẹ.

Không một ai nháy mắt, cũng không ai đưa ra hiệu lệnh. Với một lực đâm chất chứa bao năm khao khát, tủi hổ và thương nhớ của những đứa trẻ từng mất đi tất cả ở thế giới song song, Kaigou và Suzuno lao tới. Họ chồm lấy bố mẹ tôi, ôm chầm lấy họ bằng tất cả sức lực còn lại.

``\dots Mẹ?'' tiếng Kaigou run rẩy, bờ vai cô gái vốn nổi tiếng sắt đá và có thể đấm lún tường gạch giờ đây rung lên bần bật. ``Mẹ ơi\dots\ con về rồi\dots\ con đã tìm thấy mẹ rồi\dots''

``\dots Bố?'' Suzuno nghẹn ngào, ba chiếc chuông vàng trên cổ tay em rung lên những tiếng leng keng dồn dập, hòa lẫn vào tiếng khóc nấc. ``Bố ơi, đừng bỏ bọn con nữa\dots\ đừng bao giờ bỏ chúng con một mình nữa\dots''

Mẹ tôi vốn còn đang ngẩn ngơ trước cảnh tượng đột ngột, thì bàn tay bà vô thức dang ra, vòng ôm lấy hai thân hình nhỏ bé đang đổ vào lòng. Bà không hiểu chuyện gì đang xảy ra, nhưng giọt nước mắt nóng hổi của Kaigou thấm qua lớp áo len đã đánh thức bản năng người mẹ sâu thẳm trong bà. ``Ồ\dots\ các con\dots\ được rồi, không sao cả\dots'' bà thì thào, bàn tay run rẩy vuốt ve mái tóc đuôi ngựa của Kaigou, ``Mẹ ở đây rồi, không ai bỏ các con một mình nữa\dots''

Bố tôi cũng đứng sững người, chiếc kính lão ông vừa tháo ra còn nằm trong bàn tay. Ông nhìn hai cô gái lạ mặt đang khóc nức nở trong vòng tay mẹ, mắt ông cũng dần đỏ lên. Không hiểu sao, nhìn dáng vẻ đau thương của hai đứa trẻ xa lạ, trái tim ông như bị bóp nghẹt, như thể chính ông cũng đã từng trải qua khoảnh khắc chia ly này trong một giấc mơ xa xôi nào đó không có thật. ``Hai đứa\dots\ bình tĩnh nào\dots'' ông cất giọng khàn, khó khăn lắm mới nói nên lời vì chưa bao giờ gặp trường hợp nào như vậy.

``Hả?'' tôi vô thức thốt lên. Mắt tôi trợn tròn, đầu óc đột nhiên bị cuốn vào một mớ hỗn loạn không thể giải thích nổi. Làm sao mà hai cô gái này lại có thể gọi bố mẹ tôi là ``bố'' và ``mẹ'' một cách tự nhiên đến thế? Tôi nhìn chằm chằm vào lưng của Kaigou đang run lên vì khóc, rồi quay sang nhìn Suzuno, người cũng đang vùi mặt vào vai mẹ tôi mà nức nở. Câu hỏi cứ lặp đi lặp lại trong đầu tôi như một bản ghi bị kẹt: họ quen bố mẹ tôi từ khi nào? Trước đây chúng tôi có bao giờ gặp nhau đâu? 

Tôi biết hai chị em đã từng nói rằng họ đã từng có bố mẹ, và họ cũng đã qua đời cách đây bảy năm. Vậy không lẽ\dots\ bố mẹ của họ nhìn khá giống với bố mẹ tôi? Tôi đứng yên tại chỗ, không thể cử động, chỉ có thể nhìn cảnh tượng trước mắt mà không thể nào lý giải được.

Tôi đứng đơ như cây cơ tại chỗ, não bộ hoàn toàn đóng đóng băng, không tài nào xử lý nổi lượng thông tin vừa bùng nổ trước mắt. Vừa lúc đó, từ gầm giường nơi thường ngủ quên, con chó Mi trắng nhỏ nhắn hơn mười ba tuổi của nhà tôi bật dậy, lắc đuôi mạnh rồi sủa mất tiếng ``gâu gâu'', như thể đang nghi ngờ hỏi ``Ai đấy, cậu chủ? Sao trong nhà có thêm nhiều người lạ thế?''. Nó không dám lại gần, chỉ đứng ở xa vừa sủa vừa vẫy đuôi, nhận ra có gì đó lạ nhưng vẫn cảm nhận được hơi ấm thân quen từ những người mới đến.

Đúng lúc đó, Ryuuhou - lúc này mới về tới nhà sau buổi tập ban nhạc, lưng còn đeo bao đàn guitar - bước qua cánh cửa. Cô bé nhìn cảnh tượng hỗn loạn trước mắt, cặp sách trên tay suýt rơi xuống đất.

Em tôi tiến lại gần, ghé sát tai tôi thì thầm bằng giọng sững sờ: ``\vie{\hl{Anh ơi, họ\dots{} họ là ai vậy? Sao cô bé có chỏm tóc ngố kia\dots{} lại có gương mặt giống em thế? Và\dots{} cô bé tóc tết này là ai vậy?}}''

Tôi không đáp nổi một lời. Trước mắt tôi, mẹ tôi dù chưa hiểu chuyện gì đang xảy ra nhưng bản năng người mẹ đã khiến bà siết chặt hai cô gái lạ mặt vào lòng. Bàn tay ấm áp của bà vuốt ve mái tóc họ một cách tự nhiên, dịu dàng như thể bà đã từng làm điều ấy hàng nghìn lần trong tiền kiếp. Bố tôi đứng im lặng, gỡ chiếc kính lão ra, đôi mắt ông ánh lên thứ ánh sáng sâu thẫm của một người vừa gặp lại đoạn ký ức tưởng chừng đã thất lạc ở một chân trời xa xăm. Bố mẹ hướng mắt về phía tôi, như muốn hỏi rằng ``hai đứa này là ai vậy con?'' và ``chuyện gì xảy ra với họ vậy?''

Đứng bên cạnh tôi, Naomi khẽ nắm chặt lấy vạt áo phông của tôi. Đôi mắt xanh lục bảo của thiên thần nhỏ rưng rưng, im lặng nhìn cảnh đoàn tụ gia đình ấm áp - thứ cảm giác tình thân trọn vẹn mà linh hồn tồn tại 1300 năm của em chưa từng được trải qua.

\lang{Nhật}

Ở góc phòng khách, Coco và Aloe đang ngó đầu ra từ sau cánh cửa, hai đôi mắt mở to hết cỡ trước màn kịch kịch tính không kém gì phim truyền hình dài tập.

``Kuon-paisen dẫn bạn về nhà sao?'' Coco giọng có mang chút hiếu kỳ, miệng còn đang nhai nửa miếng bánh mochi mà cô vừa lấy từ tủ lạnh, dường như chẳng bận tâm gì tới bầu không khí của gia đình. 

Aloe đứng bên cạnh cô thì toát mồ hôi hột, giọng nhỏ hết sức có thể: ``Kuon-kun\dots\ Anh rước thêm ai về nữa vậy? Tại sao họ lại nhận bố mẹ anh là bố mẹ họ thế?''

``\dots Anh cũng đang ngơ ngác như bò đội nón như các em đây này,'' tôi thốt lên khi hai đứa em bị bắt quả tang đang rình xem cảnh tượng kia, vừa thở dài vừa xoa nhẹ đầu Naomi vẫn đang đứng nắm chặt áo tôi, ``Anh cũng không ngờ mọi chuyện lại trở nên phức tạp nhanh đến thế này đâu.''

\ct

Sau khi xoa dịu những cảm xúc bùng nổ ban đầu, cả nhà quây quần bên bàn ăn phòng khách. Bầu không khí vẫn đọng lại sự ngỡ ngàng, lúng túng nhưng tràn ngập sự tò mò.

Bố tôi rót một ấm trà nóng, đặt từng tách trà xuống trước mặt ba vị khách lạ ấy. Ông ôn tồn cất giọng: ``Nào, bây giờ các cháu có thể cho cô chú biết tên của các cháu được không? Và các cháu tới từ đâu nào?''

Kaigou khẽ cúi đầu, bàn tay cô siết chặt lấy vạt áo, cất giọng trầm lắng: ``Cháu tên là Kaigou\dots\ Hikawa Kaigou ạ.''

Suzuno tiếp lời, đôi mắt hai màu lấp lánh sự lễ phép: ``Còn cháu là Suzuno\dots\ Hikawa Suzuno. Rất vui được gặp chú cô ạ.''

Cô bé nhỏ tuổi nhất ngồi giữa hai người quay sang nhìn bố mẹ tôi, đôi mắt xanh lục bảo xoe tròn ngây thơ: ``Cháu là Naomi\dots\ Cháu chào chú, cháu chào cô ạ.''

Bố mẹ tôi ngẩn người khi nghe thấy cái họ ``Hikawa'' cất lên từ miệng hai cô gái. Cả Coco và Ryuuhou sững lại khi nghe tới cái tên đó, em tôi còn để rơi cả gói bim bim đang ôm trên tay: ``H-Hikawa!? Sao hai chị cũng mang họ giống nhà mình vậy?''

Nàng Rồng thì nhanh nhảu bổ sung ngay: ``Trời đất, không lẽ ngoài Kuon-paisen và Ryuuhou-chan, nhà mình còn có thêm hai đứa em lạc ở vũ trụ song song nào đó ư!?''

Mẹ tôi đặt tách trà xuống, ngập ngừng hỏi: ``Kaigou\dots\ Suzuno\dots\ Tại sao vừa nhìn thấy chúng ta, hai cháu lại khóc và gọi là bố mẹ?''

Kaigou hít một hơi thật sâu để nén lại cơn nghẹn ngào. Cô lấy từ trong túi áo ra một chiếc điện thoại thông minh khá cũ kỹ, màn hình đã xuất hiện vài vết nứt. Cô thao tác mở một tấm ảnh chụp gia đình lưu trong bộ nhớ máy rồi xoay màn hình hướng về phía bố mẹ tôi.

``Bởi vì\dots\ cô chú trông giống hệt bố mẹ quá cố của chúng cháu,'' Kaigou ngước nhìn bố mẹ tôi, giọng cô run lên. ``Từ ánh mắt, nụ cười, cho tới cả cách nói chuyện trầm ấm của chú\dots\ Tấm ảnh này được chụp từ mười năm trước, khi gia đình cháu vẫn còn đầy đủ.''

``Khi còn đầy đủ?'' cả bố mẹ tôi cùng lúc lên tiếng. ``Ý của cháu là\dots''

Suzuno khẽ gật đầu như ngầm xác nhận điều mà ông bà định nói: ``\dots Đúng là như vậy ạ. Bố mẹ cháu đã qua đời cách đây bảy năm rồi. Mẹ cháu thì gặp nạn trong lúc giao hàng, còn bố cháu thì cũng không qua khỏi vì tai nạn lao động ạ.''

Bầu không khí ở trong phòng bỗng chốc nặng trĩu, lặng đi chỉ còn tiếng thở nhẹ của những người trong phòng. Aloe sau khi nghe câu nói ấy, liền bật lên giọng run rẩy đầy thương cảm: ``Trời ơi\dots\ thế ra từ bảy năm trước chị em mình đã phải một mình chống chọi với mọi thứ sao? Tội nghiệp ghê\dots''

Suzuno liền ngước lên, đôi mắt hai màu lấp lánh nhưng miệng vẫn nở một nụ cười nhỏ để an ủi mọi người: ``Không sao đâu ạ\dots\ bọn en đã vượt qua hết rồi. Chị không cần phải buồn nữa đâu.''

Mọi người trong phòng khách cùng ghé đầu vào nhìn màn hình điện thoại. Trong ảnh là một người đàn ông trung niên mang gương mặt, chòm râu và vết sẹo nhỏ trên lông mày trái y hệt bố tôi, bên cạnh là một người phụ nữ có nụ cười hiền hậu, ánh mắt dịu dàng không khác gì mẹ tôi. Hai đứa trẻ trong ảnh chính là Kaigou và Suzuno thời thơ ấu.

Ryuuhou thốt lên ngạc nhiên: ``Trời ơi\dots\ Y như đúc! Giống như thể được chụp từ cùng một người vậy!''

Mẹ tôi nghe xong câu chuyện thì lặng người, đôi mắt đỏ nhòe vì thương cảm. Bà nhìn Kaigou và Suzuno, rồi lại nhìn sang cái họ ``Hikawa'' mà cả hai đang mang từ thế giới cũ.

``Thế ra\dots'' Bố tôi trầm ngâm, ngón tay gõ nhẹ lên mặt bàn gỗ. ``Ở thế giới của các cháu, các cháu cũng mang họ Hikawa\dots\ và cũng có một người bố, người mẹ giống hệt chúng ta sao?''

Kaigou cúi đầu, xưng hô một cách lễ phép nhưng đong đầy sự kính trọng: ``Vâng\dots\ thưa chú. Ở thế giới của cháu, bố mẹ đã mất từ lâu\dots\ Khi nhìn thấy cô chú, cháu\dots\ cháu đã không kiềm chế được.''

Suzuno khẽ gật đầu, đôi mắt hai màu lấp lánh: ``Chúng cháu xin lỗi vì đã đường đột như vậy ạ\dots''

Coco khoanh hai tay dựa lưng vào khung cửa phòng, chớp chớp mắt đảo quanh một lượt từ tôi sang Kaigou, rồi lại đăm đăm ngắm nghía Suzuno và em gái Ryuuhou. Bất chợt, em ấy đập tay cái *BỐP* vào đùi cái mạnh, khiến cả mấy người trong phòng đều giật mình: ``Này này, có ai để ý không? Kuon-paisen với Kaigou-san này đâu phải chỉ giống nhau một chút đâu, họ chính là cặp song sinh thật sự chứ còn gì nữa! Từ đường nét gương mặt, cho đến cái ánh mắt cứ như thiếu ngủ triền miên ấy, trùng khớp hết mấy!''

Ryuuhou nghe thế liền há hốc mồm, liên tục chồm người lên để so sánh gương mặt tôi với Kaigou. Em vừa chỉ vào trán tôi rồi lại chỉ vào trán Kaigou, thốt lên kinh ngạc: ``Ôi trời ơi! Cái nếp nhăn nhỏ xuất hiện mỗi khi anh trầm tư cũng y hệt nhau luôn! Mà mái tóc xoăn tự nhiên của anh, chị Kaigou cũng có luôn, trông như in ra vậy!''

Bố tôi ngồi đối diện, vừa điều chỉnh gọng kính vừa đưa mắt nhìn qua nhìn lại mấy lần giữa hai đứa, cũng gật gù lắc đầu ngỡ ngàng: ``Đúng thật là\dots\ Nếu không nói trước, ai mà nghĩ hai đứa không phải anh em song sinh. Cái vẻ ngoài cứ như được nhân bản ra ấy.''

Mẹ tôi thì gật đầu lia lịa, hai tay vẫn cầm chiếc khăn giấy lau vội nước mắt, miệng cười tươi rói: ``Phải đó! Con trai mẹ có mái tóc xoăn giống hệt ông ngoại, Kaigou cũng na ná như thế luôn! Trời ơi, đúng là máu mủ ruột thịt thì làm sao mà tách được, tìm nhau mãi thôi.''

Kaigou từ lúc nào cũng chết trân nhìn tôi, hai má bỗng ửng hồng lên như đào mùa xuân khi nghe mọi người đều trầm trồ như vậy. Cô khẽ nắm chặt lấy vạt áo, giọng nói nhỏ như muỗi kêu: ``Thật\dots\ thật sự giống vậy ạ cô chú?''

Suzuno bên cạnh cũng hào hứng gật đầu tán thành, đôi mắt hai màu lấp lánh niềm vui: ``Từ lúc gặp Onii-sama ở The Void, em đã cảm thấy quen thuộc lắm! Em đã bảo rồi mà, hai người đúng là anh em ruột thịt mà!''

Ryuuhou cũng thêm vào, liên tục vỗ tay thích thú: ``Thảo nào lúc Kaigou-san mới bước vào nhà, em đã cảm thấy có một sự gắn kết lạ! Giờ mới rõ, hóa ra là anh em ruột mà!''

``Đâu có!'' cả tôi và Kaigou cùng lúc đồng thanh phản đối, giọng nói đồng đều đến mức cả phòng bật cười. Trời ạ, sao đến cả bố mẹ và em gái tôi cũng đồng tình với mấy cô bạn vậy?

\dots Cũng chẳng trách được ai. Hai đứa có gương mặt giống nhau đến thế, họ không nói là song sinh cũng mới là lạ.

Aloe đứng cạnh Coco cũng gật gù tán thành hết mực: ``Đúng á đúng á! Còn Suzuno-chan với Ryuuhou-chan nhà mình thì cũng là hai bản sao song sinh luôn! Chỉ khác là Suzuno-chan có chỏm tóc ngố, và màu mắt hai cô bé thì đảo ngược vị trí thôi, còn lại thì y hệt đúc!''

Nàng Rồng nhếch môi cười hếch gáy, ánh mắt đảo xuống phần ngực của hai cô gái mười lăm tuổi rồi bĩu môi: ``Đâu chỉ có thế! Nhìn vòng một mà xem! Suzuno mới mười lăm tuổi mà đã đẫy đà, trù phú phơi phới, trong khi Ryuuhou nhà mình thì\dots\ ôi giời ơi, lép kẹp như cái sân bay kia kìa!''

``Cái gì cơ!?'' Ryuuhou nghe vậy thì mặt đỏ bừng vì tức giận. Cô bé nhăn mặt, khoanh tay trước ngực ngoảnh mặt đi nơi khác, hừ một tiếng rõ to: ``Kệ em! Liên quan gì tới chị chứ Coco-san!''

Cả phòng khách bật cười trước phản ứng trẻ con của Ryuuhou. Giữa bầu không khí náo nhiệt ấy, Naomi ngồi lặng lẽ bên cạnh Suzuno. Cô bé tỏa ra một aura vô cùng dễ thương, nhẹ bẫng và trong lành như làn gió xuân. Dù không nói nhiều, nhưng sự hiện diện dịu dàng của thiên thần nhỏ tám tuổi như xoa dịu mọi căng thẳng, khiến lòng ai cũng thấy bình yên kỳ lạ.

Kaigou đặt nhẹ tách trà xuống đĩa, ngón tay cô vô tình miết vào lớp gỗ mài mòn của mặt bàn, như thể đang tìm lại cảm giác thực tại sau chuỗi ngày dài vô định. Không khí trong phòng bỗng chìm vào sự im lặng cô đọng, mọi người đều hồi hộp chờ đợi câu chuyện đầu tiên được hé mở.

Cô ngước nhìn bố mẹ tôi, giọng trầm thấp nhưng bình thản: ``Cháu và Suzu\dots\ chúng cháu không thuộc về thế giới này. Ở thế giới song song của cháu, mọi thứ bắt đầu sụp đổ khi một vết nứt thứ nguyên tỏa ra ánh sáng đỏ rực nuốt chửng tất cả. Khi bị cuốn vào Hư Vô (The Void), cháu và em gái đã bị chia cắt, rơi vào khoảng không trắng xóa vô tận không lối thoát. Nếu không có Kuon-kun nghe thấy tiếng gọi và kiên trì rạch phá ranh giới để kéo từng người chúng cháu ra, giúp hai chị em đoàn tụ, có lẽ chúng cháu đã mãi mãi tan biến trong hư vô rồi.''

Tôi gật đầu, nhớ lại khoảnh khắc đầu tiên gặp hai người ở miền không gian hỗn loạn, nơi những mảnh vỡ của các thế giới va đập vào nhau liên tục. Lúc đó tôi - một người đi lại đầy lạc lõng ở trong thế giới hầu như không có một hướng chỉ dẫn gì và súyt bị chính cô ấy đập chết - nghe thấy dòng năng lượng yếu ớt của Suzuno truyền qua màn không gian. Sau khi thuyết phục được Suzuno rằng tôi không phải người xấu, hai chúng tôi mới dò theo mỏ neo ký ức để tìm ra Kaigou đang kiệt sức ở một góc rạn nứt khác. Cảm giác giúp hai chị em ôm chầm lấy nhau giữa The Void thực sự rất nghẹn ngào.

``Vậy thế giới của hai chị em trông như thế nào vậy?'' Aloe khẽ thì thầm, đôi mắt xanh nhũ của em tròn xoe đầy tò mò sau khi nghe Kaigou kể về cảnh tan hoang của thế giới song song.

``Ở thế giới của bọn em, tỷ lệ nam và nữ bị lệch hoàn toàn, chỉ còn 1 nam cho 3 nữ,'' Suzuno tiếp lời chị mình, giọng em run rẩy khi kể lại những ngày tháng kinh hoàng. ``Đàn ông ở đó coi phụ nữ như đồ vật, họ bắt cóc những cô gái trẻ để bán cho người giàu, hoặc ép làm việc nặng nhọc suốt ngày đêm. Bố mẹ thì mất sớm, chỉ có hai chị em nương tựa vào nhau trong một thế giới ngày càng lụi tàn và lạnh lẽo. Từ đó, chỉ có Onee-sama và em nương tựa vào nhau, chạy trốn khắp nơi để không bị bắt. Cũng vì vậy mà em đã học chậm mất hai năm.''

Mẹ tôi hơi sững lại, vừa định rót thêm trà vào tách của Suzuno thì chậm tay một nhịp. ``Học chậm hai năm ư? Vì không có tiền đóng học phí sao?'' bà thốt lên, giọng ngậm ngùi như vừa nghe một điều gì đó quá đỗi thương tâm.

Bố tôi ngồi bên cạnh cũng gật gù lắc đầu, thở dài: ``Hai đứa đúng là\dots\ Ryuuhou nhà mình cũng bị `chậm tiến độ' vì toàn mải mê với nhóm nhạc mà bỏ bê bài vở, cũng tận hai năm đây.''

Nghe bị bố nhắc tên, Ryuuhou mặt đỏ bừng, liền hắt hơi đẩy nhẹ vai bố: ``Bố ơi!'' Coco thì hừ mũi vỗ đùi cười lớn: ``Nghe chưa Ryuuhou-chan? Người ta học chậm vì chạy trốn sinh tử, còn em thì học chậm vì toàn chạy trốn đi tập bài tập về nhà! Giống nhau cái nỗi gì đâu mà nói chung chung!'' 

Em tôi bị réo tên đúng vào chỗ yếu, liền ôm lấy cổ Coco lắc lư để trút giận: ``Coco-san lại nói xấu em nữa! Em đã cố gắng học tập nhiều rồi mà!'' Cả phòng bật cười trước cảnh nhốn nháo của hai đứa, bầu không khí vốn nặng nề dần dịu lại, trước khi tôi tiếp tục câu chuyện.

Tôi thở dài, nhớ lại lần đầu tiên Suzuno kể với tôi về quá khứ của mình, khi chúng tôi đang ngồi nghỉ chân trên một hòn đá vụn vặt. Lúc đó em ấy đã ngồi khóc nức nở trong vòng tay chị mình, nói rằng em sợ rằng mình sẽ lại bị rơi vào tay những kẻ xấu. ``Suzuno đã từng suýt bị hãm hại ngay ở trong nhà bên thế giới của họ. May mắn là Kaigou đã phát hiện kịp thời và đấm tơi bời hắn ta. Em ấy cũng đã từng chịu đựng bị các bạn nam bắt nạt ở trường nữa, nên hầu như chẳng có chỗ nào gọi là an toàn cho hai chị em cả.''

Kaigou siết chặt nắm đấm, ánh mắt bùng lên sự kiên định: "Bởi vì thế giới cũ quá tàn nhẫn, cháu đã thề sẽ dùng toàn bộ sức mạnh của mình để bảo vệ Suzuno và Naomi!"

Vừa nói, Kaigou vừa vô tình chống tay xuống mặt bàn gỗ sồi dày dặn của nhà Hikawa để biểu thị quyết tâm.

*RẮC!*

Một tiếng động khô khốc vang lên. Chiếc bàn gỗ sồi nguyên khối khẽ rung lên, góc bàn nơi tay Kaigou tì xuống xuất hiện một vết nứt kéo dài. Cả nhà Hikawa chúng tôi - từ bố mẹ cho tới Ryuuhou, Coco và Aloe - đồng loạt đứng hình, mồ hôi hột đổ ra như tắm trước sức mạnh ``quái vật'' có thể nhấc bổng cả xe tải của cô gái tóc đuôi ngựa.  

``S-Sức mạnh khủng quá\dots'' bố nuốt nước bọt, gượng cười khẽ dịch tách trà ra xa.

Tôi nhìn khung cảnh đó mà chỉ có thể gãi đầu cười trừ: ``À\dots\ con quên dặn, ở bên kia cô ấy từng đấm lún tường gạch đấy ạ\dots\ Thậm chí còn nâng được cả một chiếc xe ben nặng chục tấn mà không xi nhê gì cơ.''

``Cái gì cơ?!'' Ryuuhou mặt xám đi, chỉ tay vào vết nứt đang lan rộng trên mặt bàn, ``Chị ấy thật sự làm được vậy ạ!?'' Coco cũng há hốc mồm không nói nên lời, vừa định đập tay lên bàn để tỏ vẻ ngạc nhiên thì nhớ ra chuyện vừa xảy ra, tay liền dừng lại giữa không khí, gãi cổ mà không dám đặt xuống. Aloe đứng cạnh cô cũng lắc đầu liên tục, mắt cứ dán chặt vào bàn tay của Kaigou như sợ cô ấy lại vô tình làm vỡ thêm gì nữa.

Suzuno vội vàng giải thích trước khi mọi người còn tưởng chị mình là quái vật: ``Chuyện là\dots\ hồi ở thế giới cũ, Onee-sama bị một trận sốt siêu vi lạ hành hạ gần chết. Bác sĩ bảo không qua khỏi được, nhưng sau một tuần nằm liệt giường, chị ấy bỗng dưng khỏe lại hoàn toàn. Kể từ đó sức lực của chị ấy bỗng chốc như vậy, đến nay thì không ai đụng độ được.''

Naomi, lúc này đang nắm chặt vạt áo Suzuno, ngước mắt to tròn nhìn mọi người rồi nhỏ giọt kể thêm như nhớ ra chuyện gì: ``Naomi cũng nghe Kaigou-oneechan kể hồi ở nơi em sống, có một đám con trai đến bắt nạt Suzuno-oneechan, chị ấy đã đánh chúng đến nhừ tử mà không một ai dám lại gần can ngăn đâu ạ! Bọn chúng sợ chị ấy lắm, sau đó không ai dám đến tìm hai chị em gây chuyện nữa.'' Tôi cá là em ấy đang nhắc tới chuyện hồi lúc chúng tôi làm nghi lễ ở Aogami, Kaigou đã một mình ``thanh toán'' đám con trai đã bắt cóc Suzuno rồi làm nhục em ấy. Cái đó thì quả thực họ xứng đáng bị như vậy.

Mẹ tôi buông chiếc khăn giấy đang cầm trong tay, bước lại gần ngồi cạnh Kaigou, tay nhẹ nhàng đặt lên mu bàn tay của cô gái như sợ cô ấy sẽ tự làm tổn thương mình: ``Vậy\dots\ cuộc sống thường ngày có vất vả không cháu? Phải luôn kiềm chế sức mạnh để không làm vỡ đồ đạc, hay vô tình làm tổn thương người khác, chắc khó khăn lắm nhỉ?''

Cô nàng tóc đuôi ngựa ngước nhìn mẹ, ánh mắt sắc lạnh ban đầu bỗng mềm lại, cô khẽ gật đầu thành thật: ``Dạ\dots\ có những lúc rất khó. Cháu phải luôn nhớ phải giảm lực khi cầm đồ, bắt tay người khác, sợ rằng sẽ vô tình làm vỡ hay làm đau ai. Trước đây ở thế giới cũ, cháu đã từng làm gãy mấy cái xương của mấy tên bắt cóc khi tự vệ, cũng từng làm vỡ cửa nhà nhiều lần vì không kiểm soát được lực\dots''

``Nghe có vẻ cực khổ nhỉ,'' Aloe nghe tới đây mà đổ mồ hôi hột.

Tôi ngồi bên cạnh lắng nghe, trong đầu bắt đầu nghĩ ngợi: phải có cách nào để giúp cô ấy. Game có thể thiết kế một loại vòng đeo tay tích hợp bộ điều biến lực, có thể giới hạn sức mạnh của cô ấy trong phạm vi an toàn khi ở nhà, để cô ấy không phải lúc nào cũng căng thẳng kiềm chế nữa. Hoặc thậm chí là một phòng tập riêng có tường và sàn được gia cố bằng vật liệu đặc biệt, để Kaigou có thể thoải mái giải phóng sức mạnh mà không sợ phá vỡ thứ gì. Tôi sẽ bàn với ông ta vào tối nay để vạch ra kế hoạch cụ thể.

Trong lúc tôi còn đang mải nghĩ cách, Aloe ngồi yên lặng lắng nghe câu chuyện, hai bàn tay em nắm chặt chẽ trên đầu gối. Đôi mắt cô nàng rướm lệ, giọng nói khàn lại đầy phẫn nộ: ``Những tên khốn đó\dots\ sao chúng lại có thể đối xử tàn nhẫn như vậy được? Nếu có chị ở đó, chị sẽ không để chúng đụng đến các em một sợi lông tơ nào.''

Tôi vốn định nhắc đến việc bản thân cô nàng Succubus cũng từng trải qua quãng thời gian yếu đuối tương tự, nhưng cuối cùng lại kìm lại lời nói. Khoảnh khắc này không phải lúc thích hợp để nhắc đến những chuyện quá khứ ấy.

Mẹ tôi trong vô thức tự nhiên kéo Suzuno vào lòng, tay nhẹ nhàng vỗ lên lưng em như an ủi một đứa trẻ vừa thoát khỏi cơn ác mộng khủng khiếp. Bà thì thầm lặp đi lặp lại: ``Tội nghiệp hai cháu thật\dots\ tội nghiệp quá\dots'' Bố tôi cũng chẹp miệng một tiếng, vẻ mặt đang cố gắng nén lại cơn giận dữ trước những đau thương mà hai cô bé phải gánh chịu: ``Chú phải thừa nhận, hai cháu đã thực sự rất mạnh mẽ mới có thể sống sót qua được một nơi chẳng khác nào địa ngục như thế.''

Ryuuhou lúc này cũng rơi nước mắt, em vội đứng dậy chạy tới ôm chặt lấy Suzuno, như thể muốn dùng thân mình che chắn cho người chị em song sinh của mình khỏi mọi vết thương còn sót lại. ``Không sao nữa đâu Suzuno-san, bây giờ chúng ta cùng ở đây rồi, không ai có thể làm hại hai chị được nữa.''

Cũng đứng dậy ngay sau đó, Aloe ngồi xuống bên cạnh Kaigou và đặt một bàn tay ấm áp lên vai cô, giọng nói dịu dàng vỗ về: ``Đúng đó, Kaigou-san, ở đây thật sự an toàn. Từ giờ hai chị không còn phải lo lắng về bất cứ điều gì nữa đâu.''

Coco cũng hào hứng không kém, bẻ cổ tay tạo ra tiếng lách cách rõ rệt: ``Đúng đó! Chị đây sẽ là vệ sĩ riêng cho hai người! Nếu có chuyện gì không ổn thì cứ hô to tên Coco là có ngay!''

Naomi liền chầm chậm buông tay rồi nhỏ nhắn bò qua, vòng hai cánh tay nhỏ bé ôm chầm lấy cả hai chị. Đôi mắt xanh lục bảo của em ướt đẫm nước mắt, giọng nói nhỏ nghẹn lại: ``Naomi thương Kaigou-oneechan, thương Suzuno-oneechan lắm\dots\ Từ giờ Naomi sẽ cùng các chị ở đây, không ai có thể bắt nạt các chị nữa đâu.''

Kaigou cảm nhận bàn tay ấm áp của Aloe trên vai, rồi tiếng thì thầm non nớt của Naomi vang lên, bờ vai cô gái vốn đã cứng rắn trải qua bao nhiêu sóng gió bỗng rung lên bần bật. Bảy năm trôi qua trong sự cô độc, chạy trốn khắp nơi, bàm chặt hơi ấm duy nhất chỉ có từ em gái mình, giờ đây cả căn phòng này tràn ngập sự quan tâm, những vòng tay che chở mà họ tưởng đã mãi mãi mất đi cùng với bố mẹ ở thế giới cũ. 

Suzuno cũng nghẹn ngào, ngẩng đầu lên nhìn những gương mặt thân thiện xung quanh, cảm giác hơi ấm lan tỏa từ mẹ tôi đang vỗ nhẹ lưng em, như làn nước ấm thấm vào mảnh đất khô cằn, làm tan đi mọi nỗi sợ hãi đã bám riết hai chị em trong suốt bảy năm ròng rã. 

``Cảm ơn\dots\ cảm ơn mọi người,'' Kaigou khẽ nói, giọng cô run rẩy như không dám tin rằng khoảnh khắc ấm áp này lại có thể đến với mình.

\ct

Sau một lúc im lặng, khi những giọt nước mắt đã bớt rơi, Naomi mới ngước lên từ vai tôi, đôi mắt xanh lục bảo của em nhìn quanh mọi người trong phòng. 

Ryuuhou là người đầu tiên tiến lại gần, ngồi xuống bên cạnh em và cầm lấy bàn tay nhỏ bé còn đang run nhẹ của Naomi, giọng en ấy mềm đi đáng kể: ``Naomi-chan, em có chuyện gì muốn kể với mọi người không? Nếu khó quá thì không cần nói cũng được đâu ạ.'' 

Aloe cũng ngồi quỳ xuống phía đối diện, đôi mắt rướm lệ gật đầu phụ họa: ``Đúng đó, em đừng ép bản thân nếu những ký ức ấy vẫn còn khiến em đau. Nhưng nếu muốn chia sẻ, bọn chị đều ở đây để lắng nghe hết lòng.'' 

Coco vừa vỗ nhẹ vào vai em với nụ cười ấm áp hiếm thấy: ``Chúng ta đều là một gia đình rồi, có chuyện gì cũng có thể cùng nhau san sẻ, đừng giữ nó trong lòng một mình nữa ha.''

Naomi ngước nhìn những gương mặt quan tâm của mọi người, sau đó khẽ mím môi rồi bắt đầu nói, giọng nhỏ nhẹ nhưng rõ ràng: ``Naomi cũng không phải từ thế giới này. Em sống ở thời Nara, cách đây đã 1300 năm rồi. Lúc đó, làng của Naomi hay xảy ra sạt lở đất, các cụ bảo là phải hiến tế một đứa trẻ cho thần núi để xoa dịu thần núi. Naomi là người được chọn, họ đưa cháu lên ngôi đền và lấp cháu cùng với những đứa trẻ khác dưới lòng đất\dots''

Ryuuhou không còn kiềm chế được cảm xúc nữa, liền ôm chặt lấy thiên thần nhỏ vào lòng: ``Sao mà có thể tàn nhẫn đến thế được!? Một đứa bé mới chỉ tám tuổi mà lại phải gánh chịu những chuyện kinh khủng như vậy sao?''

``Tại sao mỗi câu chuyện đều bi kịch như vậy\dots?'' bố tôi bất giác thốt lên với giọng tràn đầy thương cảm, rõ ràng vừa ngán ngẩm vừa xót xa trước những số phận đau thương mà ba đứa trẻ phải chịu đựng.

Tôi nhẹ nhàng vuốt ve mái tóc mềm mại của Naomi, bỗng nhớ lại lần đầu tiên em ấy kể với tôi về quá khứ của mình - đó là lúc mà tôi và nhóm lớp 1-C lần đầu tiên đối mặt với chính nơi em ấy từng bị chôn vùi. Tôi tiếp lời câu chuyện của em: ``Naomi đã phải mắc kẹt dưới dạng một linh hồn lang thang suốt 1300 năm ở đó, cho đến khi bọn con tìm thấy em ấy. Tại nơi đó, chúng con đã dùng năng lượng Thanh tẩy để hồi sinh em thành người thực, để em có thể được sống một cuộc đời bình thường như bao đứa trẻ khác.''

Tôi vốn cũng muốn kể về việc Game đã hỗ trợ chúng tôi hoàn thành sứ mệnh ấy ra sao, nhưng cuối cùng lại thôi, không muốn làm mọi chuyện phức tạp thêm; hơn nữa ông ta cũng đang trong trạng thái nghỉ ngơi, không tiện để nói chuyện cùng mọi người. Dù sao đi nữa, ông ta vốn là người không mấy thích giao tiếp với người lạ.

Bố mẹ tôi đứng người ngẩn ngơ khi nghe hết câu chuyện của Naomi, họ không bao giờ ngờ được rằng đứa bé tám tuổi xinh xắn đang ngồi trước mặt lại phải trải qua một quá khứ đau thương đến thế. Mẹ tôi nghẹn ngào kéo Naomi vào lòng cùng với Suzuno, ôm chặt hai đứa nhỏ như thể sợ rằng chỉ cần buông tay một chút, em sẽ biến mất không còn dấu vết. ``Một đứa bé mới tám tuổi mà phải trải qua những chuyện như vậy\dots\ Naomi ngoan lắm, bây giờ con đã có bố mẹ, có anh chị em bên cạnh rồi, không ai sẽ bỏ con một mình nữa đâu.''

Bố tôi ngồi bên cạnh, bàn tay vuốt nhẹ mái tóc mềm mại của em, giọng ông trầm xuống đầy thương cảm: ``Vậy bây giờ, khi đã được sống như một đứa trẻ bình thường ở thế giới này, cháu thấy thế nào, Naomi-chan? Có quen với mọi thứ ở thời hiện đại không?''

Naomi ngước đôi mắt xanh lục bảo long lanh từ lòng mẹ, miệng nở một nụ cười trong veo như nắng mai: ``Naomi thấy mọi thứ đều kỳ diệu lắm ạ! Ánh đèn điện sáng rực như sao sa, cái ô tô bốn bánh chạy nhanh hơn ngựa hoang, còn cả cốc parfait mà Onii-chan cùng hai chị đã mua cho cháu cũng ngon lắm. Tất cả đều mới lạ và tuyệt vời quá!'' giọng nói non nớt của cô bé như một dòng suối mát lành chảy qua căn phòng, tan đi mọi nỗi niềm nặng trĩu vừa nãy. Cả nhà ai cũng mềm lòng trước sự trong trẻo ấy, như thể mọi đau thương của em chưa từng làm xước đi được sự ngây thơ thiên thần mà em mang theo.

Tôi ngồi đó, bất chợt nhớ lại khoảnh khắc định mệnh ở Aogami, khi chúng tôi vừa hoàn thành nghi thức hồi sinh cho em. Lúc đầu, Kaigou, Suzuno và tôi đã từng có ý định để em ở lại thị trấn ấy - chúng tôi nghĩ rằng giữa cảnh núi non bình yên, em sẽ được mọi người và các bạn học lớp 1-C chăm sóc tốt hơn là đi lang thang cùng bọn tôi, những kẻ luôn dính vào chuyện rắc rối. 

Nhưng Naomi đâu có chịu nghe. Em bướng bỉnh ôm chặt lấy chân váy của Suzuno, đôi mắt to tròn ngấn nước mà không chịu rời ra, lắc đầu từ chối ráo riết. Khi Ryuuhou nghe thấy ý định ích kỷ ấy - dù cho đó là ý tốt - của chúng tôi, em ấy lập tức phản đối kịch liệt, nói rằng chúng tôi không thể vứt bỏ đứa bé mới vừa có được cuộc đời mới. 

Em ấy đã nói một câu mà đến bây giờ tôi vẫn còn nhớ mãi, khi tôi phải đưa ra lựa chọn rất khó khăn cho con bé: ``Naomi muốn được ở bên cạnh mọi người. Dù chỉ thêm một giây thôi cũng được. Nếu phải tan biến\dots\ thì tan biến cùng nhau vẫn tốt hơn là sống sót một mình ở đây.'' Chính sự chân thành đó đã đánh gục tất cả chúng tôi, khiến chúng tôi không thể nào nói thêm lời nào để buộc em ở lại. 

Ryuuhou lúc đó cũng vỡ òa, lao tới ôm chặt lấy em, nói rằng em sẽ không bao giờ để em phải một mình nữa. Từ đó, chúng tôi mới thực sự là một gia đình. Em ấy nũng nịu ôm chặt lấy cô bé tám tuổi từ phía sau, cằm tựa nhẹ lên đỉnh đầu em: ``Naomi-chan đáng yêu ghê, từ giờ chúng mình sẽ là chị em tốt nhất của nhau nhé, hứa vậy!''

Coco cười lớn để xua tan bầu không khí buồn bã đang bao trùm căn phòng. ``Thôi thì những chuyện cũ đã qua hết rồi! Bây giờ ba đứa đều là thành viên chính thức của nhà Hikawa, chúng ta sẽ cùng nhau sống thật vui vẻ nhé! Ngày mai mình cùng đi mua sắm đồ dùng mới cho các chị nhé, tôi biết một cửa hàng thời trang rất đẹp ở trung tâm thành phố!''

Aloe gật đầu hớn hở liền theo: ``Đúng đó đúng đó! Tôi cũng sẽ làm bánh matcha để các chị cùng ăn vào ngày mai, chị Kaigou chắc chắn sẽ thích món đó lắm.''

Tôi ngồi yên lặng ngắm nhìn cảnh tượng trước mắt, cảm giác một sự ấm áp lạ lùng dâng tràn từ trong lòng. Ba người đã phải lang thang xuyên qua vô số thế giới, cuối cùng cũng tìm được một nơi thực sự có thể gọi là nhà.

Tôi ngước nhìn Kaigou, người vẫn đang ngồi im lặng nhưng đôi mắt cô ấy đã không còn cái sắc lạnh kinh hoàng của ngày đầu chúng ta gặp nhau, thay vào đó là sự nhẹ nhõm tột độ, như thể một tảng đá lớn đã được nhấc ra khỏi trái tim cô sau bao nhiêu năm chờ đợi.

Rồi ánh mắt tôi chuyển sang Suzuno, cô bé vốn luôn nhút nhát và chỉ dám bám chặt vào bóng dáng chị mình để tìm kiếm sự an toàn, giờ đây cũng đang nở một nụ cười ấm áp khi vòng tay của mẹ tôi ôm lấy em. Ánh sáng trong đôi mắt hai màu của em không còn cái hoảng sợ luôn trực xuất như ngày đầu lạc lõng giữa The Void, mà thay vào đó là sự bình yên, như thể cuối cùng cũng tìm được nơi em có thể buông bỏ hết mọi phòng thủ.

Và cuối cùng là Naomi, cô bé từng là một linh hồn lang thang suốt 1300 năm chỉ khao khát được giải thoát khỏi cõi giới vô định, giờ đây đang gục đầu vào vai Ryuuhou, đôi mắt xanh lục bảo long lanh niềm vui khi được mọi người vây quanh. Không còn cái vẻ lặng lẽ rụt rè ban đầu, em đã trở thành một đứa trẻ đúng với lứa tuổi của mình, được nuông chiều, được yêu thương và có một mái nhà thực sự để gọi là nhà.

Cả ba người, ba số phận từng lang thang xuyên qua bao chiều không gian để tìm kiếm một nơi chốn, cuối cùng đều có cùng một điều ước đã thành hiện thực: một mái nhà, một gia đình thực sự.

\ct

\lang{Etsunan}

Tối hôm ấy, căn phòng ngủ vốn rộng rãi của tôi bỗng trở thành nơi trú tạm cho cả năm anh em ruột, thêm một con chó già nhõng nhẽo. Một tấm đệm cỡ lớn được trải phẳng giữa sàn nhà, và năm người chúng tôi - Kuon, Kaigou, Suzuno, Ryuuhou và Naomi - đều ngồi xếp thành một hàng dài trên đó, bên cạnh là Mi, con chó trắng cảnh đã hơn mười ba năm tuổi của nhà Hikawa, đang cuộn tròn ngay giữa chân cô bé tám tuổi như đã quen thuộc từ lâu.

Tôi vừa vuốt nhẹ mái đầu nhỏ của Naomi đang ngủ gật dựa vào vai mình, bàn tay cũng vô tình xoa nhẹ bộ lông mềm của Mi-chan đang dụi vào chân tôi, thì liền bị một cái đẩy nhẹ từ phía Kaigou làm ngả người: ``Này, cậu có thể ngồi nghiêm chỉnh một chút được không? Từ nãy giờ cứ dựa qua đây làm gì, nghĩ chị em tôi là gối dựa sao?'' cô nhếch mép, giọng nói vẫn mang chút cứng đầu quen thuộc.

Ryuuhou lập tức phụ họa ngay, vừa bỏ một miếng snack vào miệng vừa lắc đầu cười nhạo, tay cũng thả một mẩu nhỏ bánh quy cho Mi đang ngước nhìn cô: ``Đúng đó anh! Từ lúc ở The Void, anh đã biết dựa dẫm chị Kaigou rồi hả? Mà cũng tài lẻ đấy, đi ra ngoài một thời gian mà về nhà lại có thêm ba cô gái, đúng là `đi xuyên không mà rước cả gia đình về'!'' em vừa nói vừa cười khúc khích, nhưng ánh mắt nhìn tôi lại ẩn chứa nỗi nhẹ nhõm không thể che giấu, rằng cuối cùng anh mình cũng trở về an toàn sau những sự kiện ấy, thậm chí còn mang theo những người bạn đồng hành để cùng nhau sống trong ngôi nhà này.

Tôi liền vỗ nhẹ vai em mình, giả vờ tức giận: ``Này em nói gì thế? Anh có phải là kẻ đi bắt cóc đâu mà gọi là rước về! Là các em ấy đồng ý ở lại nhà mình thôi đấy!''

Kaigou vừa định mở miệng châm chọc thêm câu nào, thì Suzuno đã nhẹ nhàng nắm lấy cánh tay chị mình, giọng nói nhỏ nhẹ như làn gió, tay cũng vô tình vuốt qua bộ lông trắng bông của Mi-chan đang chồm lên sưởi ấm: ``Chị ơi, đừng trêu anh Kuon nữa mà\~{}'' em vừa nói vừa nhẹ nhàng điều chỉnh chiếc chăn mỏng đang che cho Naomi, đảm bảo cô bé không bị lạnh, sau đó lại quay sang cười dịu dàng với chúng tôi: ``Mà đúng là, ở đây ấm áp thật sự, em chưa bao giờ cảm thấy an tâm như thế này kể từ khi bố mẹ ở thế giới cũ mất đi.''

Cả không gian phòng bỗng chìm vào một sự ấm áp nhẹ nhàng, như thể mọi mệt mỏi của những ngày lang thang đều tan biến đi trong ánh đèn bàn vàng ấm. Chúng tôi vừa tiếp tục cười đùa nhỏ tiếng, thì cánh cửa phòng khẽ mở ra, Bố mẹ tôi bước vào, trên tay bê theo một đĩa hoa quả đã được gọt sẵn và vài ly sữa nóng hổi. Vẻ mặt của hai người vẫn chưa hết bàng hoàng sau những biến cố lớn đã xảy ra trong ngày, phía sau lưng bố mẹ còn có Mi-chan vẫy đuôi nhỏ, nhảy lúp xúp vào phòng như cũng muốn cùng tham gia buổi tụ tập ấm áp.

Bố tôi kéo chiếc ghế đẩu ngồi xuống đối diện năm anh em. Ông nhấp một ngụm trà, chậm rãi kể lại một chi tiết khiến tôi nổi toàn bộ da gà:

``Thực ra\dots\ mấy năm gần đây, cả bố và mẹ đều thường xuyên thấy những giấc mơ kỳ lạ,'' bố tôi trầm ngâm nói. ``Trong giấc mơ, bố mẹ thấy một cô gái tóc đuôi ngựa và một cô bé có chỏm tóc ngố đứng giữa màn sương mù dày đặc. Họ cất tiếng gọi 'Mẹ', gọi 'Bố' bằng thứ giọng nói da diết mà rất lâu rồi bố mẹ không được nghe từ Ryuuhou. Còn có con bé Naomi nữa, cùng con Mi nhà mình đứng bên cạnh, trông hết sức thân thuộc, như thể đã ở cùng chúng ta cả đời rồi.''

Mẹ tôi mỉm cười, tiến lại gần nắm lấy bàn tay nhỏ nhắn của Naomi, tay nhẹ nhàng vuốt bộ lông của Mi-chan đang dụi vào chân bà: ``Nhìn con bé ấy, lòng mẹ lại thấy bình yên đến lạ kỳ. Như thể con vốn đã thuộc về ngôi nhà này từ rất lâu rồi.''

Tôi nghe mà rợn cả sống lưng: có khi nào dòng máu và mối liên kết tâm linh thực sự biết cách tìm thấy nhau, bất chấp khoảng cách của những chiều không gian? Ngay cả con Mi già vốn chỉ quấn quýt mỗi bố mẹ và Ryuuhou từ trước, bây giờ cũng nhanh chóng thân thiết với ba cô bé mới đến, chẳng hề có chút xa lạ nào.

Kaigou nhìn bố tôi, giọng cô nghẹn lại: ``Những giấc mơ đó\dots\ có lẽ là lúc cháu và Suzu gào khóc trong tuyệt vọng nhất khi bố mẹ ở thế giới bên kia qua đời\dots''

Suzuno ôm lấy cánh tay Mẹ tôi, đôi mắt hai màu lấp lánh giọt lệ, tay vẫn liên tục xoa nhẹ đầu của con chó đang nằm gọn bên cạnh: ``Cháu cảm ơn cô chú\dots\ Cảm ơn vì đã cho chúng cháu cảm giác được trở về nhà\dots''

Ryuuhou nghiêng đầu nhìn Suzuno rồi lại nhìn Kaigou, khẽ mỉm cười: ``Em không biết lý do tâm linh là gì, nhưng nhìn hai người\dots\ em thực sự cảm thấy giống như chị em trong nhà vậy. Ngay cả con Mi nhà mình cũng vậy, mới gặp có một ngày mà đã bám dính lấy cả ba chị em như thể đã biết từ lâu.''

Naomi ngước đôi mắt xanh lục bảo lên nhìn Kuon rồi lại nhìn bố mẹ, tay nhỏ bé ôm chặt lấy con chó đang nằm trong vòng tay mình: ``Naomi cảm thấy\dots\ nơi này rất ấm áp. Naomi muốn ở lại đây mãi mãi, cùng với Mi-chan nữa.''

Lời nói non nớt của cô bé như làm tan đi lớp màn ngại ngần cuối cùng còn vương vấn giữa mọi người, để căn phòng tràn ngập niềm vui ấm áp của một gia đình đang thêm đầy đủ thành viên. 

Coco vốn đứng nép bên cánh cửa từ lúc nào, chờ đúng khoảnh khắc này liền bước vào với nụ cười rạng rỡ, phá tan sự yên ấm nhẹ nhàng bằng giọng nói hào hứng: ``Ha ha! Thế là từ nay nhà Hikawa đông vui phải biết nhé! Kuon-paisen hết thời làm con một rồi! Sau này anh sẽ phải nhường nhịn nhiều thứ đấy, còn phải cạnh tranh với cả con Mi để được ôm Naomi-chan nữa!''

Tôi liền chống tay sườn lắc đầu đáp lại: ``Thế em nghĩ em gái anh Ryuuhou làm cảnh à? Dù có kiếm được nhiều tiền và trông hoạt bát thế nào, đến bữa vẫn có lúc phải đòi tôi mua parfait, hay bắt tôi sửa máy tính cho em ấy đâu, anh vốn đã là anh cả từ lâu rồi mà!''

Aloe reo lên vui vẻ gật đầu liền: ``Vậy sáng mai em sẽ phụ mọi người dọn dẹp thêm phòng ngủ luôn! Cũng phải chuẩn bị thêm một chiếc chăn nhỏ cho Mi-chan nữa, để con bé cũng có chỗ ngủ ấm áp trên tầng ba cùng các em.''

``Ừ, cảm ơn em nhiều. Tầng ba nhà mình rộng lắm, vốn là phòng mẹ và Ryuuhou đang ở, đủ không gian để Kaigou, Suzuno và Naomi cả ba chuyển vào ở chung. Còn anh thì vẫn ở phòng cũ cùng bố trên tầng năm thôi.''

Nhà tôi có một nguyên tắc bất di bất dịch từ trước đến nay là nam nữ không thể ở chung phòng ngủ, để giữ phép tắc và tránh phiền phức không đáng có. Đó là điều hiển nhiên mà gia đình đông con nào cũng thể cả mà, đúng không?

Naomi vừa nghe xong thì mím nhẹ môi, đôi mắt xanh lục bảo long lanh thoắt buồn, em nhỏ giọng bật ra, tay vẫn ôm chặt con chó trong lòng: ``Vậy là\dots\ Naomi không được ngủ chung với Onii-chan nữa sao ạ?''

Kaigou nghe vậy liền chọc ghẹo tôi ngay, nhếch mép nói: ``Ồ, hóa ra cậu cũng có lúc sợ ở chung với bọn tôi hả? Chỉ được cái nói lý do ngon lành, chứ có phải không dám ngủ chung với mấy đứa con gái bọn tôi đâu, thậm chí cả con Mi cũng tranh cướp được sự quan tâm của cậu cơ mà.'' 

Ryuuhou ở bên cạnh thì chỉ cười trừ không nói gì, thầm nghĩ bản thân em ấy cũng không ít lần lẻn qua ngủ ké ở phòng hai bố con tôi tầng năm, lúc nào cũng chẳng hề ngại ngần gì, còn thường xuyên mang theo cả con Mi nữa để cùng ngủ chung. Còn Suzuno thì chỉ nhẹ nhàng cười, em cũng chẳng mấy mặn mà với cái quy định đó, nếu có cơ hội em vẫn sẵn sàng ngủ chung với tôi mà không chút e thẹn.

Bố tôi gật đầu, ánh mắt đong đầy sự bao dung của người cha: ``Dù thế nào đi nữa, kể từ hôm nay, các cháu cứ coi đây là nhà của mình. Chúng ta sẽ cùng nhau tìm ra câu trả lời.''

Tiếng cười nhỏ lan tỏa khắp căn phòng, ánh đèn vàng ấm áp phủ lên những gương mặt rạng ngời, còn có tiếng thở đều nhẹ nhàng của con chó già đã ngủ say trong vòng tay Naomi, như thể mọi nỗi sợ hãi và lang thang trong vô tận chiều không gian bấy lâu nay, cuối cùng cũng tìm được bến đỗ bình yên thực sự. Ngôi nhà nhỏ này, từ hôm nay sẽ mãi là mái nhà chung của tất cả chúng ta.

\begin{center}
    \uline{Ngày hôm sau.\\Bệnh viện trung tâm thành phố.}
\end{center}

Cả nhà kéo nhau tới bệnh viện trung tâm thành phố để làm xét nghiệm ADN cũng thật tình cờ. 

Lý do là bởi lẽ, hôm qua Coco và Aloe đã trêu chọc mãi rằng tôi trông vô cùng giống Kaigou, còn Ryuuhou thì lại có nét tương đồng kỳ lạ với Suzuno như thể hai cặp song sinh bị lạc mất từ nhỏ. Tôi bèn đùa lại rằng nếu vậy thì đi xét nghiệm cho chắc, nếu không phải thì cũng chẳng tốn kém bao nhiêu (cùng lắm cũng chỉ là tốn thời gian), ai ngờ cả nhà lại đồng loạt hào hứng đòi đi thật. 

Hơn nữa, hôm nay tôi cũng được A-san thông báo nghỉ làm, công việc dự án chưa quá gấp để phải vội vàng, lịch học ở trường cũng không có buổi nào bận rộn. Ryuuhou thì hôm qua đã hô hào xin nghỉ học một ngày với lý do ``gia đình có việc quan trọng'', mà tôi biết thừa đó chỉ là cái cớ để trốn tránh bài kiểm tra môn toán sắp tới, đúng như những gì Coco đã chỉ điểm trước đó hôm qua. Thế là cả nhà kéo nhau tới bệnh viện trung tâm thành phố để thực hiện xét nghiệm ADN.

Tám giờ sáng, khi chúng tôi bước vào sảnh chờ chính của bệnh viện, không khí bệnh viện ban sáng ngột ngạt mùi cồn y tế, tiếng máy rít đều đều từ các phòng xét nghiệm và tiếng người rôm rả ngoài hành lang. Ánh nắng sớm lọt qua những ô cửa kính lớn, nhưng không đủ sức làm tan đi cái hơi lạnh lẽo đặc trưng của nơi này, chỉ làm nổi bật những gương mặt vội vã chạy qua lại giữa hành lang. 

Tôi đứng giữa dòng người, lòng bỗng dấy lên một chút ngậm ngùi nhớ lại ba năm về trước, cũng tại căn bệnh viện này, tôi đã phải vật lộn với cơn đau dạ dày cấp tính trong đúng giai đoạn kỳ thi vào đại học cận kề. Hồi đó, ngày ngày tôi nằm trên giường bệnh vừa truyền nước, vừa ôm tập ghi chú ôn bài, mồ hôi lạnh túa ra vì cơn đau quặn từng cơn mà vẫn không dám bỏ phí một phút nào. Giờ quay lại đây, mọi ký ức ấy vẫn còn rõ mồn một như mới xảy ra hôm qua, và vẫn khiến tôi khẽ rùng mình khi nhớ lại cảm giác ấy.

Bên cạnh tôi, Naomi lần đầu tiên đặt chân vào một cơ sở y tế hiện đại, đôi mắt xanh lục bảo tròn xoe nhìn mọi thứ với sự tò mò pha chút ngại ngần. Cô bé chỉ vào một chiếc máy đo huyết áp đặt ở góc sảnh, kéo nhẹ tay áo tôi rồi hỏi nhỏ: ``Anh ơi, cái máy kia màu trắng kia là gì thế ạ? Nó sáng đèn nhấp nháy kìa!''

Tôi nhẹ nhàng đáp lời: ``Đó là máy đo huyết áp đấy, Naomi, để xem tim của mọi người có khỏe không.'' Cô bé gật gù như hiểu rồi, lại chỉ vào một y tá đang cầm ống tiêm đi qua, giọng nói có chút run rẩy: ``Cái que nhỏ chị kia cầm là gì ạ? Nó có làm đau không?'' 

Lần này Suzuno ngồi bên cạnh nhẹ nhàng vuốt mái tóc em, giải thích dịu dàng: ``Đó là ống tiêm để lấy một ít máu thôi, chỉ hơi chích một chút thôi, không đau nhiều đâu.''

Ryuuhou ngồi cạnh tôi cũng thở dài, tự bao giờ em ấy cũng chìm vào suy tư: ``Lâu lắm rồi mới lại vào bệnh viện. Lần cuối cùng em tới đây là hồi lớp ba, bị sốt virus phải nhập viện năm ngày liền, nằm đến mức chán cả người. Từ đó em quyết định sẽ tập thể dục nhiều hơn để không bao giờ phải quay lại đây nữa.'' Em ấy nói rồi lại lắc đầu cười, vẻ hơi ngại ngùng trước ký ức non nớt đó.

Phía đối diện, Kaigou và Suzuno lại im lặng nhìn khung cảnh xung quanh, ánh mắt phảng phất nỗi choáng ngợp trước sự đông đúc và hiện đại của nơi này. Tôi biết trong lòng hai chị em đang dậy lên những ký ức buồn ở thế giới cũ: hồi đó, cũng tại một bệnh viện tương tự, họ đã phải nhận lại thi thể của bố và mẹ sau những tai nạ không may. Ngày hôm đó, hai em đứng trước căn phòng lạnh giá của nhà xác, không dám tin rằng người thân yêu nhất của mình đã rời đi mãi mãi. Nhưng giờ đây, khi cả nhà Hikawa ngồi đầy đủ bên cạnh, những bóng ma quá khứ dần tan biến, thay vào đó là một hy vọng mới về một mái ấm thực sự.

``Được rồi mọi người ơi!'' Ryuuhou đột nhiên đứng dậy bật dậy từ ghế nhựa cứng của sảnh chờ, bàn tay đập nhẹ vào đùi mình để thu hút sự chú ý, giọng cô ấy vang lên trong không gian ồn ào như một đội trưởng đang hô hào tập hợp đội hình trước trận đấu: ``Nhiệm vụ hôm nay rất đơn giản: chúng ta sẽ đi lấy mẫu, nộp hồ sơ và chờ kết quả. Ai cũng phải làm theo đúng chỉ dẫn, không được tản mạnh nghe chưa! Đặc biệt là Naomi, em không được chạy lung tung quanh các phòng bệnh nhé, nếu muốn đi đâu phải có anh chị đi cùng!''

Em ấy giơ cao bàn tay như đang cầm một cây dùi cui chỉ huy, vai ưỡn ra với dáng vẻ dũng mãnh mà chẳng có vẻ gì là nghiêm túc, khiến vài người ngồi chờ xung quanh không khỏi quay sang nhìn chúng tôi với ánh mắt tò mò. Naomi ngồi cạnh tôi nghe vậy liền giơ hai bàn tay nhỏ bé che miệng cười khúc khích, đôi mắt xanh lục bảo nhíu lại thành hai vầng trăng non.

``Em nói nghe như chúng ta đi đánh trận vậy\dots'' tôi nhún vai cười ngán ngẩm với tinh thần nhiệt huyết\dots\ không đúng chỗ của em gái mình, vừa nói vừa vỗ nhẹ lên đầu Naomi để em bé khỏi cười quá mức. ``Chẳng phải chỉ là đi lấy máu làm xét nghiệm thôi sao, em làm lớn chuyện thế ai mà hiểu được.''

Kaigou lập tức hưởng ứng, vỗ tay to vào vai Ryuuhou đến mức cô em gái suýt ngã nhào về phía trước, tiếng cười lớn của cô ấy vang lên làm mấy người ngồi gần nữa phải ngoái đầu: ``Đúng vậy! Chị sẽ phụ trách hậu phương, canh chừng từng người một, ai dám bỏ đi đâu là chị biết liền! Dám lảng đi mua đồ ăn vặt trước khi làm xong thủ tục thì đừng trách chị!'' 

Tôi ngồi đó chỉ có thể thở dài ngán ngẩm, nhìn hai người đang hùa nhau như sắp dẫn cả đội đi thực hiện một nhiệm vụ bí mật cấp quốc gia chứ không phải đi làm xét nghiệm ADN thường ngày.

Suzuno ngồi cạnh mẹ tôi thì nhẹ nhàng lấy khăn giấy lau đi giọt nước mắt vừa rơi ra vì cười quá mức, em chỉ có thể gật đầu liên tục đồng tình với hai người, bàn tay nắm chặt lấy vạt áo mẹ như sợ rằng nếu không cẩn thận thì hai người kia thực sự sẽ kéo cả nhà chạy lung tung trong bệnh viện. Còn Naomi thì cứ lắc đầu cười không ngớt, em bé reo lên: ``Naomi không chạy đâu ạ! Em sẽ ở cạnh anh Kuon luôn!''

Bố tôi đặt lon cà phê vừa mua từ máy bán hàng tự động xuống ghế, mỉm cười nhìn hai cô gái đang hào hứng vô cùng, ông lắc đầu nói với mẹ tôi: ``Naom lúc nào cũng năng lượng dồi dào thế. Ngày nào cũng như có pin lắp vào người, không bao giờ thấy mệt mỏi.'' Bố vừa nói vừa rút điện thoại ra kiểm tra lại lịch hẹn mà chúng tôi đã đặt từ hôm qua, đảm bảo không có gì sai sót trong thủ tục.

Mẹ tôi thêm vào, tay bà nhẹ nhàng vuốt qua mái tóc mềm mại của Suzuno ngồi cạnh, giọng nói ấm áp lan tỏa: ``Đông đủ cả nhà thế này mới là gia đình chứ. Nhà mình đã lâu rồi chưa có lúc nào ồn ào và vui vẻ như thế này, từ lúc Ryuuhou lớn lên cũng ít khi thấy em ấy hào hứng đi đâu cùng cả nhà như vậy.'' 

Chúng tôi cùng nhau đứng dậy, cầm lấy tất cả các túi xách và đồ đạc cá nhân, bước theo hướng quầy tiếp tân để hoàn tất các thủ tục, cả nhà cùng hướng tới nhiệm vụ tưởng chừng đơn giản mà lại mang ý nghĩa thay đổi cả cuộc đời.

Từng bước chân chúng tôi vượt qua hành lang bệnh viện đầy người, lướt qua những y tá cầm khay thuốc chạy vội, những bệnh nhân ngồi chờ khám với vẻ mặt mệt mỏi, nhưng trong lòng mỗi người lại tràn ngập một niềm vui nho nhỏ, một sự mong chờ khó tả. Không ai trong chúng tôi thực sự nghĩ rằng kết quả xét nghiệm sẽ có gì bất thường, mọi người chỉ coi đây là một trò đùa vui, một chuyến đi chơi cuối tuần cùng gia đình sau thời gian dài mỗi người bận rộn với công việc và học tập. Nhưng sâu thẳm trong mỗi trái tim, có một chút hồi hộp vô hình, một chút mong mỏi rằng điều kỳ diệu nào đó sẽ xảy ra, rằng những người bạn đồng hành mà chúng tôi tình cờ gặp được trên con đường lang thang, thực sự sẽ trở thành một phần máu mủ của gia đình này.

Tất cả như đang đứng trước một cánh cửa bí ẩn, không ai biết phía sau đó là gì, nhưng tất cả đều sẵn sàng bước qua cùng nhau, chuẩn bị cho một bản án, một lời tuyên bố sẽ thay đổi toàn bộ cuộc đời và số phận của mỗi người trong chúng tôi mãi mãi.

\ct

Hai tiếng chờ đợi trôi qua dài như hai kỷ nguyên. Chúng tôi lần lượt đi lấy mẫu máu: Naomi khẽ run tay khi kim tiêm chạm vào da, nhưng vẫn cắn môi không kêu la, chỉ nắm chặt lấy ngón tay tôi cho đến khi y tá rút kim ra mới thở phào nhẹ nhõm. Suzuno cũng ngồi sát bên mẹ, đôi mắt hai màu dời đi chỗ khác khi lấy máu, còn Kaigou thì bình thản như thể chẳng có gì xảy ra, dù tôi thấy móng tay cô ấy bấm nhẹ vào lòng bàn tay đến trắng bệch.

Chúng tôi quay lại ghế chờ, nán lại với những câu chuyện vu vơ để xoa tan hồi hộp, thấy Ryuuhou lôi điện thoại ra chụp ảnh chung với Naomi. Nhưng chẳng ai thực sự tập trung vào những lời nói đó, ánh mắt cứ nhổn nhổn hướng về phía cửa phòng xét nghiệm, chờ đợi điều gì đó mà chúng tôi cũng chưa dám gọi tên.

Cuối cùng, khi vị bác sĩ trưởng khoa bước ra với xấp hồ sơ trên tay, vẻ mặt ông đẫn ra như vừa nhìn thấy hiện tượng siêu nhiên. Ông đẩy cặp kính lên mũi, lướt qua từng người trong chúng tôi với ánh mắt khó tin, miệng há hốc muốn nói mà mãi mới thốt ra được: ``Gia đình Hikawa\dots\ đây là kết quả của mọi người. Tôi làm ở khoa xét nghiệm này đã hơn ba mươi năm mà chưa bao giờ thấy trường hợp nào kỳ lạ như thế này.''

Tôi đứng dậy đầu tiên, bước tới nhận tập hồ sơ từ tay bác sĩ, lòng bàn tay toát mồ hôi lạnh. Cả nhà xúm lại sau lưng tôi, hơi thở của mọi người trở nên gấp gáp khi tôi lần lượt lật từng trang giấy trắng tinh, những con số và chữ phân tích gen hiện ra trước mắt dần dần làm rõ một sự thật không thể tin được.

Không, không phải vì tôi hay Ryuuhou là con nuôi như mấy bộ manga drama mì ăn liền mà tôi thường đọc trên mạng. Hai anh em tôi thực sự đúng là anh em ruột, cái đó thì miễn bàn, giấy trắng mực đen rõ ràng. Tôi chỉ vào dòng đầu tiên, nơi xác nhận quan hệ huyết thống giữa tôi và Ryuuhou, như thể để tự thuyết phục bản thân rằng đó là sự thật.

Nhưng xuống tới những dòng dưới, chúng tôi mới thực sự bị choáng ngợp: mẫu ADN của Kaigou, Suzuno và Naomi với bố mẹ, với tôi và Ryuuhou, tất cả đều trùng khớp đến 99,99\%, tương đương với quan hệ ruột thịt trực hệ!

``\eng{\dots Thế này là thế nào đây, Game?}'' tôi vô thức gõ nhẹ vào mặt đồng hồ, cất giọng run rẩy hỏi, dường như vẫn chưa thể tin được vào mắt mình.

Giọng nói quen thuộc của ông ta khẽ hiện lên trong tâm trí tôi, phát ra một tiếng bíp nhỏ đủ để chỉ tôi nghe thấy, như thể sẵn sàng giải đáp mọi câu hỏi đang dâng trào trong đầu: ``\eng{Sự trùng khớp ADN này là hệ quả tự nhiên của các sự kiện đã diễn ra. Khi các thực thể vượt qua lỗ xoáy thứ nguyên kết nối thế giới, cấu trúc không-thời gian tại điểm giao thoa đã tạo ra một sự cộng hưởng năng lượng đặc biệt. Năng lượng này, kết hợp với hiệu ứng từ nghi thức Thanh tẩy đã thực hiện trước đó, đã tự động đồng hóa và tái cấu trúc hệ gen của ba cô gái để phù hợp với dòng máu của gia đình cậu đấy.}''

``\eng{Vậy\dots\ còn Naomi thì sao?}''

``\eng{Đối với thực thể Naomi, mặc dù ban đầu là một linh hồn tái sinh không mang vật chất, quá trình cộng hưởng gen đã diễn ra liên tục khi cô bé gắn bó mật thiết với Suzuno trong suốt chuyến du hành xuyên các chiều không gian. Khi quá trình tái cấu trúc hoàn tất, hệ gen của Naomi cũng hoàn toàn trùng khớp với hệ gen của Suzuno, và theo tính chất bắc cầu sinh học được thiết lập, cô bé chính thức mang trong mình dòng máu của gia đình Hikawa, không khác gì các anh chị em ruột thịt khác trong nhà.}''

Nghe những lời ông ta giải thích xong tôi đứng chết trân giữa phòng, đầu óc vẫn đang cố gắng tiếp nhận toàn bộ thông tin quá sức tưởng tượng ấy. Game lại tiếp tục hiển thị một bảng phân tích chi tiết, như để làm rõ hơn mọi thứ:

\begin{itemize}
    \item Hikawa Kaigou: Em gái song sinh của Hikawa Kuon, sinh sau Kuon đúng 10 phút theo ghi nhận sinh học sau tái cấu trúc, chính thức trở thành chị cả của các em gái trong gia đình.
    \item Hikawa Suzuno: Chị em song sinh của Hikawa Ryuuhou, sinh trước Ryuuhou 5 phút, cùng mang dòng máu trực hệ của bố mẹ Hikawa.
    \item Hikawa Naomi: Con gái út của gia đình, 8 tuổi, hệ gen hoàn toàn phù hợp với quan hệ ruột thịt, chính thức mang họ Hikawa sau quá trình cộng hưởng và tái cấu trúc.
\end{itemize}

Cả thế giới dường như đảo lộn hoàn toàn trước mắt tôi. Những cá nhân từng lang thang vô định xuyên qua vô số chiều không gian, những người bạn đồng hành mà tôi tình cờ gặp được trên con đường đầy gian khổ, bỗng chốc trở thành máu mủ ruột thịt, những người thân thích nhất trên đời. Những cái tên Kaigou, Suzuno, Naomi từng chỉ hiện lên trong những trang nhật ký ghi chép về các cuộc phiêu lưu xuyên thứ nguyên, giờ đây đã chính thức có thể được ghi bổ sung vào sổ hộ khẩu gia đình, trở thành một phần không thể tách rời của mái nhà Hikawa. Mọi chuyện xảy ra nhanh đến mức tôi không kịp chuẩn bị tâm lý, không kịp nghĩ xem mình sẽ đối mặt với tất cả những thay đổi này như thế nào.

Cứ thế, suốt mấy giây dài cả căn phòng xét nghiệm lặng ngắt như tờ, chỉ còn tiếng thở gấp gáp của mọi người hòa vào tiếng máy điều hòa chạy rì rào từ góc phòng. Bàn tay tôi lạnh ngắt, từng ngón tay run rẩy quét qua từng dòng chữ in trên giấy kết quả, như thể muốn kiểm tra xem liệu có phải mình đang nhìn nhầm, hay đây chỉ là một giấc mơ hoang đường nào đó sau một đêm mất ngủ.

Mẹ tôi là người đầu tiên phá vỡ sự im lặng. Bà bước tới, tay run rẩy cầm lấy tờ giấy kết quả từ tay tôi, nước mắt không kìm được rơi lã chã lên mặt giấy trắng tinh, làm nhòe đi vài dòng chữ in. Bà ôm chầm lấy Suzuno và Naomi đứng cạnh, giọng nói nghẹn ngào xen lẫn niềm vui sướng tột độ: ``Trời ơi\dots\ trời ơi có thật không này\dots\ Tôi có thêm con gái rồi\dots\ tôi có thêm ba đứa con gái nữa rồi\dots'' Bà lần lượt vuốt ve mái đầu của Kaigou đang đứng cạnh tôi, của Suzuno, của Naomi, như thể sợ rằng chỉ cần chạm nhẹ thôi, tất cả sẽ tan biến như một giấc mơ đẹp.

Bố tôi thì đứng yên một chỗ, tay chậm rãi gỡ chiếc kính lão ra khỏi mắt, ông lấy vạt áo sơ mi lau đi lau lại đi chiếc kính, miệng cứ lẩm bẩm mãi một câu vì quá khó tin, không thể chấp nhận được sự thật kỳ diệu vừa xảy ra. Sau mấy giây, ông mới ngẩng đầu lên, nhìn cả đàn con đang đứng xúm xít trước mặt, giọng nói run rẩy nhưng tràn đầy tình yêu thương: ``Vậy là\dots\ ta có thêm ba đứa vịt giời nữa rồi.''

Naomi ngước nhìn tôi với đôi mắt xanh tròn xoe, kéo nhẹ áo tôi hỏi: ``Anh Kuon, vậy là\dots\ bây giờ Naomi cũng là em gái của anh đúng không ạ? Cũng là con của bố mẹ Hikawa ư?'' 

Suzuno thì đứng bên cạnh mẹ, nước mắt đã bắt đầu tuôn rươi, cô bé gật đầu liên tục như không thể tin được, còn Kaigou thì đứng yên lặng, tay nắm chặt lấy vai áo tôi đến nhăn nhúm, đôi mắt sắc lạnh thường ngày bỗng dưng long lanh ngấn nước. Ryuuhou thì lao tới ôm chầm lấy Suzuno, la lên trong nước mắt: ``Suzuno-san! Vậy là chúng ta thực sự là chị em ruột rồi à? Liệu\dots\ mình có nằm mơ không đây\dots?''

Tôi đứng giữa phòng xét nghiệm, giữa những tiếng khóc xúc động của mẹ, những tiếng cười nghẹn ngào của các em, cảm giác như vừa bị một lực vô hình đẩy tọt vào vai trò ``anh cả'' trong một vở kịch gia đình quy mô lớn mà bản thân chưa từng được đọc kịch bản. Từ trước đến nay, tôi chỉ có mỗi em gái Ryuuhou để chăm sóc, để nhường nhịn, bỗng chốc tôi có thêm ba người em gái nữa, ba người mà tôi đã đồng hành qua bao nguy hiểm, bây giờ lại chính thức trở thành máu mủ ruột thịt của mình. Cảm giác vừa lạ lẫm, vừa ấm áp lạ thường tràn dâng trong lòng, như thể một khoảng trống vô hình nào đó trong tim tôi bỗng được lấp đầy, từ nay trở đi không còn trống trải nữa.

Tôi quay sang Kaigou đang đứng cạnh. Tờ giấy kết quả xét nghiệm trong tay cô ấy nhăn nhúm lại từng chút một, những ngón tay dài vốn luôn giữ được sự vững vàng trong mọi tình huống nguy cấp nhất, giờ đây lại run rẩy không kém gì một đứa trẻ sắp mất đi nơi nương tựa. Đôi mắt sắc lạnh thường ngày, vốn đã xem qua biết bao cảnh tượng khắc nghiệt của thế giới cũ mà không hề nao núng, giờ đây đong đầy những giọt nước mắt trong trẻo, chỉ chực trào ra khỏi đường viền mắt. 

Cô ấy từ từ quay người, để cả gương mặt đối diện với tôi, đôi môi vốn luôn sắc sảo và cứng rắn khẽ rung lên như sắp vỡ ra: ``\dots\ Aniki (兄貴).'' Cách gọi ``anh trai'' mang đầy sự kính trọng, vốn chỉ thường xuất hiện giữa những người anh em máu mủ trong giới xã hội đen, giờ đây lại thốt ra từ miệng cô ấy một cách tự nhiên như thể đã được ấp ủ trong lòng hàng trăm, hàng ngàn ngày.

Tôi đứng ngẩn người, đầu óc vẫn chưa kịp xử lý hết những thông tin vừa nhận được, nên lúng túng quay sang nhìn cô ấy: ``Aniki? Cô đang nói tới ai--``

Chưa kịp để tôi nói hết câu, cô ấy đã bước nhanh tới, hai tay dang rộng lao tới ôm chầm lấy tôi bằng một lực siết mạnh đến mức tôi cảm thấy xương sườn mình vừa bị ép chặt, thở hơi nào cũng trở nên khó khăn (ơn giời cô ấy vẫn còn tỉnh táo để tiết chế được sức mạnh siêu phàm của mình, nếu không thì có lẽ lồng ngực của tôi đã nát bét từ lúc nào rồi).

Cô ấy gục đầu vào ngực tôi, mái tóc dài hơi thơm mùi dầu gội hoa oải hương mà mẹ vừa mua cho các con gái từ hôm qua, che đi gương mặt đang rơi lệ của cô ấy. Bờ vai cao ráo của Kaigou rung lên từng đợt, từng đợt một, kèm theo những tiếng khóc nghẹn ngào như muốn trút bỏ hết tất cả những gánh nặng đã đè lên vai cô ấy suốt bảy năm qua. 

Từ lúc bố mẹ qua đời ở thế giới cũ, cô ấy đã phải vừa làm chị, vừa làm bố làm mẹ để nuôi nấng Suzuno, phải đối mặt với những kẻ săn đuổi, phải tìm cách kiếm ăn qua ngày, phải gồng mình chống chọi với mọi cơn bão tố của cuộc đời mà không bao giờ được phép yếu đuối. Giờ đây, cuối cùng cô ấy cũng tìm được một điểm tựa an toàn, một người anh trai lớn để có thể dựa vào, để thôi không phải lúc nào cũng giữ vững bức tường thành lạnh lùng bảo vệ mình.

Suzuno đứng bên cạnh, tay vẫn ôm chặt Naomi vào lòng, mỉm cười dịu dàng khi nhìn cảnh tượng trước mặt, nước mắt cũng rơi trên má, nhưng đó là những giọt nước mắt của niềm vui. Ryuuhou đứng cạnh tôi cũng lau vội nước mắt bằng gấu áo, em vỗ nhẹ vào lưng Kaigou đang vẫn gục trên ngực tôi, tiếng nói run rẩy: ``Em vẫn không thể tin là chị Kaigou là song sinh của anh Kuon luôn đấy\dots\ Chúc mừng chị nhé.'' 

Tôi nghiêng đầu hỏi Suzuno, người vẫn đang mỉm cười nhìn cảnh tượng vừa diễn ra: ``Nhưng mà Suzuno này, tại sao cô ấy lại gọi anh là Aniki cả thế nhỉ? Nghe lạ quá.''

Cô em gái của cô nàng nhẹ nhàng nghiêng đầu, tiếp tục câu chuyện của mình để giải thích cho bố mẹ tôi còn đang ngẩn ngơ đứng nhìn: ``Ở thế giới bên kia, chị Kaigou luôn phải gánh vác mọi thứ một mình để bảo vệ em. Chị ấy chưa bao giờ cho ai thấy mình yếu đuối, chưa bao giờ khóc trước mặt chúng em, vì chị nói phải làm tấm gương vững chắc để chúng em có thể tin tưởng.

Em ấy dừng một lúc lâu, ngước nhìn cô chị vẫn đang đứng lúng túng bên cạnh tôi, rồi nhỏ nhẹ tiếp tục: ``Nhưng thỉnh thoảng trong đêm, khi em tỉnh dậy đi vệ sinh, có nhiều lúc em bắt gặp chị ấy ngồi một mình ở cửa sổ, lặng lẽ khóc vì nhớ bố mẹ, vì sợ rằng mình không đủ mạnh để bảo vệ hai chị em. Chị ấy luôn khao khát có một người anh trai lớn để tựa vào, có người để chia sẻ những gánh nặng trên vai\dots''

Suzuno nhìn tôi một hồi, tay khẽ gạt những dòng lệ vui mừng trên khóe má. ``Và qua tất cả những gì anh Kuon đã thể hiện suốt chuyến đi xuyên các chiều không gian, từ lúc anh một mình đối mặt với những hiện tượng siêu nhiên để bảo vệ chúng em, đến lúc anh luôn sẵn sàng chia sẻ miếng ăn cuối cùng, lo lắng cho chúng em từng li từng tí, anh chính là người anh trai hoàn hảo nhất mà chị ấy tìm kiếm bấy lâu nay. Em cũng vậy\dots\ em cũng đã coi anh là anh trai từ lâu rồi, từ khoảnh khắc mà anh đã giúp em đoàn tụ với chị ở The Void đó. Anh còn nhớ không?''

Tôi ngẫm lại, đúng là khoảnh khắc ở The Void đó, khi tôi đã len lỏi trong thế giới trống rỗng ấy để cô bé có thể ôm chầm lấy chị Kaigou của mình, tất cả vẫn như mới xảy ra hôm qua. Từng tiếng khóc nghẹn của em khi ấy, hơi ấm run rẩy trong vòng tay, tất cả đã khắc sâu vào tim tôi như một lời thề không cần nói ra: rằng tôi sẽ luôn bảo vệ các em, sẽ không bao giờ để bất cứ ai phải chia lìa nữa.

Naomi, vẫn đang ngậm ngùi ngồi trong vòng tay của chị Suzuno, cũng nhẹ nhàng gật đầu đồng tình, đôi mắt xanh lục bảo long lanh nước mắt mà miệng vẫn nở nụ cười ngọt ngào: ``Naomi cũng gọi anh Kuon là Onii-chan! Em thích được anh bế, thích anh kể chuyện cổ tích trước khi ngủ, thích anh mua parfait dâu tây cho Naomi ăn. Bây giờ Naomi thật sự có anh trai rồi, đúng không ạ?''

Trước câu hỏi thẳng thắn của Naomi, cổ họng tôi bỗng dưng tắc nghẽn. Tôi chỉ có thể gật đầu khó nhọc, sau đó đưa tay xoa nhẹ mái tóc mềm mại của em bé khi cô ấy ngẩng lên một lần nữa, khóe mắt tự nhiên ẩm ướt vì sự ấm áp dâng tràn từ tim. Tôi từng tưởng mình sẽ đi hết cuộc đời chỉ với một em gái, còn những người bạn đồng hành chỉ là khách qua đường, ai ngờ số phận lại tặng món quà kỳ diệu này. Giờ đây tất cả chúng ta thực sự trở thành một gia đình, gắn kết bằng máu thịt chứ không chỉ những chuyến phiêu lưu chung. Giờ đây, một chân trời mới lại mở ra trước mắt tôi, một cuộc sống gia đình mà tôi chưa bao giờ dám mơ tới, đầy đủ tiếng cười, sự ấm áp và những người thân yêu luôn ở bên cạnh.

Tôi siết chặt vòng tay ôm lấy Kaigou một lần nữa trước khi cô ấy lùi ra, vỗ nhẹ vào lưng cô ấy như an ủi, giọng tôi cũng nghẹn lại khi nói với mọi người: ``Đúng vậy, từ nay có gì chúng ta cùng chia sẻ. Không ai phải gánh vác một mình nữa.'' 

Nhưng bất ngờ nào đâu có kết thúc ở đó.

Đúng lúc bầu không khí đang ngập tràn sự xúc động, cái đồng hồ đeo tay của tôi bỗng rung nhẹ một cái, và giọng nói quen thuộc của Game vang lên trong đầu, với cái tông giọng mỉa mai vốn có của ông ta, như thể đã chờ đợi từ lâu để tung ra một quả bom: ``\eng{Kuon-user, cậu còn một thông tin quan trọng chưa nghe đâu. Đừng có mà quá vui mừng đến mức quên hết mọi thứ.}'' 

Tôi vừa toan mở miệng hỏi còn chuyện gì nữa, thì một màn hình ảo tí hon đột nhiên hiện ra trước mắt, dòng chữ trên đó rõ mồn một đến mức khiến tôi muốn té ngửa ngay tại chỗ. Giọng nói của ông ta lại tiếp tục vang lên, thậm chí còn mang chút thích thú trước bộ dạng kinh ngạc không thể che giấu của tôi: ``Cậu biết không, mặc dù cậu và Kaigou-user có ADN trùng khớp hoàn toàn như anh em song sinh, nhưng chuyện tái cấu trúc gen từ biến động thứ nguyên mà tôi đã nhắc đến hồi nãy ấy, đâu chỉ đơn giản có vậy. Cô ấy sở hữu 100\% gen trội biểu hiện độc lập, theo thuật ngữ di truyền học của thế giới mà cậu đang sống, thì nó có nghĩa là\dots\ sẽ không có bất kỳ rủi ro di truyền nào nếu hai người có con với nhau.''

Tôi tròn mắt ngạc nhiên, đầu óc có chút quay cuồng trước thông tin hoàn toàn ngoài dự đoán, miệng lẩm bẩm không dám tin: ``V-Vậy có nghĩa là\dots''

``Nói một cách đơn giản nhất, ngoài đời thực thì hai người hoàn toàn có thể kết hôn một cách hợp pháp, nếu cậu thực sự muốn. Thiên thời, địa lợi, nhân hòa đều đã hội tụ đủ cả rồi đấy.''

Tôi suýt chút nữa thì sặc nước bọt chết sấy ngay tại chỗ. Lòng bàn tay tôi toát mồ hôi lạnh, mặt tôi bừng nóng lên như vừa bị đốt lửa, mắt tôi tròn xoe nhìn vào cái đồng hồ đeo tay của mình, không thể tin được vào những gì vừa nghe. Ông ta chắc chắn đang nói đùa, chắc chắn là có cái gì đó sai sót ở đây, làm sao có thể có chuyện vô lý như thế được?

``\textbf{Cái gì cơ!?}'' Ryuuhou nghe được những gì mà em ấy lảng vảng nghe được, liền gào lên, tiếng nói của em ấy vang vọng cả căn phòng xét nghiệm, làm vài y tá đi qua ngoài cửa cũng phải ngoái đầu nhìn vào. Em ấy đứng thẳng người, mồm há hốc như vừa nhìn thấy ma: ``Anh nói thật á? Lại còn có vụ đó thật á!?''

Kaigou nghe xong thì mặt đỏ bừng tới tận mang tai, từ má lên trán đều chuyển sang màu đỏ rực, cô ấy đứng im lặng vài giây, sau đó tung ra ``ánh mắt đóng băng'' sắc lạnh nhìn chằm chằm vào khoảng không khí nơi Game đang ẩn nấp, như thể muốn đóng băng cả cái đồng hồ đeo tay của tôi để ông ta không bao giờ nói thêm một lời nào. Bàn tay cô ấy nắm chặt lại, giọng nói lắp bắp nhưng đầy giận dữ: ``Ông ta đang nói cái gì thế không biết\dots'' 

Bố tôi thì ho sặc sụa, vừa uống lon nước vừa bị nghẹn, phải vỗ vào ngực mình mấy cái mới thở được, còn mẹ tôi thì bật cười khúc khắc, tay che miệng để không cười quá to, trước sự hỗn loạn đầy đáng yêu của đám trẻ. Cả căn phòng vốn đang ngập tràn nước mắt xúc động, bỗng chốc trở thành một mớ hỗn độn vui nhộn, tiếng cười nói vang vọng lên, làm tan đi mọi nỗi lo lắng còn sót lại.

Kaigou ngẩng đầu khỏi ngực tôi, đôi mắt còn đỏ hoe nhưng ánh sáng trong đó đã rực rỡ lạ thường, rồi lúng túng buông tôi ra, lùi về một bước nhỏ, tay vội lau đi những giọt nước mắt còn sót trên má như thể vừa xấu hổ vì đã thể hiện sự yếu đuối trước mặt mọi người. ``Không phải như anh nghĩ đâu Aniki! Chỉ là\dots\ chỉ là quá bất ngờ thôi!'' 

Cô ấy bật ra nhanh như bắn, nhưng má lại từ từ bừng đỏ lên khi nói câu đó, ánh mắt liếc qua tôi rồi lại vội chuyển sang nhìn chân mình, như thể sợ bị tôi phát hiện ra sự thật ẩn trong ánh mắt.

Tôi đứng đó nhìn cô ấy, bất giác mỉm cười, vốn hiểu rõ cô ấy hơn ai hết, tôi nhận ra rằng dưới lớp hung hăng chối bỏ đó, không hề có chút khó chịu nào, ngược lại, có một chút gì đó nhẹ nhàng, như thể cái ý tưởng cô ấy đang cố gắng phủ nhận, lại không hề làm cô ấy thấy khó chịu.

Naomi cũng nhẹ nhàng gật đầu đồng tình. Đôi mắt xanh lục bảo của em bé long lanh nước mắt, nhưng khi vừa nghe Game nói câu gì đó không đầu không cuối về việc kết hôn hồi nãy (dù em ấy chưa hiểu rõ ý nghĩa hoàn toàn, nhưng cũng bắt chước được bầu không khí xấu hổ đang lan tỏa), má em cũng bừng đỏ hồng, miệng vẫn nở nụ cười ngọt ngào vụng về. Suzuno vừa vỗ nhẹ lưng em bé để an ủi, vừa mỉm cười với tôi, má cô ấy cũng hơi ửng đỏ - có vẻ như cái thông tin đột ngột của Game vẫn đang làm dậy lên sự ngượng ngùng trong cả ba cô em gái.

Tôi đứng giữa căn phòng, vừa xấu hổ vừa buồn cười, nhìn Kaigou đang nổi giận, nhìn Ryuuhou vẫn chưa hết kinh ngạc, nhìn bố mẹ đang cười vui vẻ, và trong lòng tôi dậy lên một cảm giác bình yên vô cùng.

Đây chính là gia đình, một mái nhà đầy đủ tiếng cười, những trò đùa kỳ quặc và những người luôn yêu thương nhau, bất kể quá khứ đã từng như thế nào. Tất cả những chuyến lang thang xuyên các chiều không gian, tất cả những khó khăn và nguy hiểm mà chúng ta đã trải qua, cuối cùng cũng dẫn đến khoảnh khắc này, khi tất cả chúng ta thực sự có một mái nhà để quay về, có những người thân yêu để cùng chia sẻ cuộc đời. Không còn phải lo lắng về những kẻ săn đuổi, không còn phải lo lắng về việc mai này sẽ đi đâu, bởi bây giờ chúng ta đã ở nhà, và chúng ta sẽ luôn ở bên nhau, mãi mãi.

\ct

Thủ tục nhập hộ khẩu cho ba thành viên mới tưởng chừng sẽ là một cơn ác mộng hành chính, nhưng rốt cuộc lại diễn ra theo cái cách kỳ quặc và vô lý nhất có thể.

Cùng ngày vào buổi chiều, cả gia đình Hikawa chúng tôi cùng nhau di chuyển đến Ủy ban Nhân dân phường, bởi đây là thủ tục bắt buộc để hoàn tất việc nhập hộ khẩu cho ba thành viên mới của gia đình - Kaigou, Suzuno và Naomi phải chính thức được ghi nhận là con cái hợp pháp của bố mẹ Hikawa, và giấy chứng nhận xét nghiệm ADN kỳ diệu mà chúng tôi nhận được từ bệnh viện chỉ là bước khởi đầu. 

Naomi ngồi trên ghế sau xe máy của bố, đôi mắt xanh to tròn cứ nháo nhác ngắm nhìn mọi thứ dọc đường, và khi bước chân vào không gian văn phòng ủy ban với những dãy bàn gỗ xếp ngay ngắn, những tấm biển thông tin dán trên tường và tiếng máy in rít lên liên tục, em bé nắm chặt lấy tay áo tôi kéo nhẹ, liên tục chỉ vào từng vật lạ lẫm mà hỏi: ``Anh Kuon ơi, cái đó là cái gì vậy ạ? Cái cô chú đang điền vào tờ giấy trắng kia là giấy gì ạ? Sao phòng này rộng thế mà lại nhiều người ngồi làm việc cùng nhau quá?'' Tôi chỉ có thể vừa đi vừa cười xòa, lần lượt giải thích từng câu hỏi non nớt của em, trong khi cả nhà vẫn tiến sâu vào quầy làm việc hộ tịch để bắt đầu các thủ tục còn lại.

Vị cán bộ hộ tịch tròn mắt nhìn tờ kết quả xét nghiệm ADN 99,99\% mà tôi đưa ra, rồi lại nhìn ba cô gái lạ quắc từ trên trời rơi xuống, vẻ mặt bàng hoàng như sắp đình công tới nơi: ``Không có giấy khai sinh gốc thì làm sao chú làm thủ tục nhập khẩu cho các cháu được?''

Tôi vừa toan mở miệng nghĩ cách xoay sở, bỗng Kaigou đứng cạnh khẽ giật mình, tay đập nhẹ vào trán như vừa nhớ ra chuyện trọng đại: ``À đúng rồi, cháu có mang theo giấy tờ kia mà!''

Cô nàng vội lục tung trong chiếc túi xách đã sờn cũ - vật duy nhất cô mang theo từ thế giới bên kia. Từ ngăn kéo khóa bí mật sâu bên trong, cô lấy ra một chiếc phong bì nhựa bọc kín lớp chống ẩm, đưa ra trước mặt mọi người. Bên trong là hai tờ giấy khai sinh đã ngả màu thời gian, gấp cẩn thận không một nếp nhăn.

Tôi nghiêng đầu thắc mắc: ``Sao em tự dưng nhớ ra thế? Có mà giấu kỹ thế này sao không nói sớm?''

Kaigou lúng túng xoa nhẹ mép phong bì, mắt liếc sang chỗ khác: ``Từ khi qua thế giới này lúc nào cũng chỉ nghĩ đến chuyện an toàn cho các em, quên hẳn là mình còn giữ nó\dots\ Đây là di vật cuối cùng mẹ để lại cho em ở thế giới cũ, lúc nào cũng đem theo như báu vật thôi.'' Ở thế giới cũ đầy rẫy bất công, hai tờ giấy này là tấm lá chắn duy nhất chứng minh hai chị em cô là những công dân hợp pháp chứ không phải những kẻ lang thang vô gia cư. Cô luôn giữ nó bên mình như một mỏ neo tinh thần để không bao giờ quên mất nguồn cội.

Tôi cười khổ nhìn cô em gái song sinh vừa nhớ ra chuyện quan trọng: ``\dots Tiện lợi thật. Không có cái này thì chắc chúng ta phải về tay trắng hôm nay.''

Bố tôi bước tới, tay nhận lấy hai tờ giấy từ Kaigou, rồi lôi luôn giấy khai sinh của tôi và Ryuuhou từ trong cặp ra để đặt cạnh nhau. Mẹ tôi cũng xúm lại, đôi mắt bà dần mở to khi lần lượt đọc từng dòng chữ trên cả bốn tờ giấy: ``Trời ơi\dots\ thật vậy sao? Ngày tháng sinh đều trùng khớp hết cả?'' Bà chỉ vào dòng ngày sinh trên giấy Kaigou, rồi sang của tôi, sững sờ tới nỗi không nói được thêm gì.

Trên hai tờ giấy khai sinh mà Kaigou trao cho vị cán bộ, thông tin được ghi lại một cách rõ ràng: Hikawa Kaigou: Sinh vào lúc 11:10, ngày 17/07/1996. Hikawa Suzuno: Sinh vào lúc 08:20, ngày 04/04/2006.

Đối chiếu với giấy khai sinh của tôi và Ryuuhou, ngày giờ sinh của chúng tôi lần lượt là 11:00 ngày 17/07/1996 và 08:25 ngày 04/04/2006.

Ryuuhou đứng cạnh tôi cũng kêu lên kinh ngạc: ``Chị Kaigou sinh sau anh Kuon đúng mười phút! Còn chị Suzuno thì sinh trước tôi năm phút\dots\ thật sự là chị em song sinh cả hai cặp luôn!'' Naomi đứng ôm chân chị Suzuno cũng nhô đầu ra, miệng cười toe toét vỗ tay: ``Oa, giống như trong truyện cổ tích vậy ạ!'' Suzuno chỉ đứng đó mỉm cười, đôi mắt hai màu long lanh nhìn bốn tờ giấy đặt cạnh nhau, như không dám tin rằng số phận lại kỳ diệu đến thế.

``Trời đất\dots'' vị cán bộ hộ tịch soi đi soi lại mốc giờ in trên giấy rồi đưa đánh giá: ``Đúng mốc giờ ngày xưa mẹ con đi đẻ rồi! Không lệch một phút nào!''

Nhưng vừa khi mọi chuyện tưởng đã ổn thỏa, một nút thắt khó gỡ lại xuất hiện: rắc rối lớn nhất lại nằm ở cô bé út của chúng ta, Naomi. Em hoàn toàn trắng tay, không hề có bất kỳ giấy tờ tùy thân nào từ quá khứ. Làm sao chúng ta dám nói với ai rằng cô bé vốn là một linh hồn xuất thân từ mười ba thế kỷ trước, trải qua hàng trăm năm lang thang mới được tái sinh vào thể chất hiện tại? Nếu mở miệng ra nói vậy, chắc chắn vị cán bộ phường sẽ gọi điện đến bệnh viện tâm thần để đưa cả nhà đi kiểm tra mất.

Tôi đứng đó bối rối, liếc nhìn Kaigou đang cũng nhíu mày suy nghĩ, thì nghe em gái song sinh của tôi khẽ thì thầm vào tai: ``Em quên thật rồi, Naomi không có giấy tờ gì cả\dots\ Bây giờ phải làm sao đây?''

Tôi vừa toan vỗ vai cô ấy để an ủi, thì Naomi bỗng kéo nhẹ gấu áo tôi, đôi mắt xanh to tròn đầy lo lắng: ``Anh Kuon ơi\dots\ Naomi có làm gì sai không ạ? Sao chú bác lại nhìn Naomi với vẻ khó hiểu thế?''

Suzuno nhanh chóng bước tới ôm chầm lấy em út, xoa nhẹ mái tóc mềm mại của Naomi để em bình tĩnh, rồi cũng ngước nhìn tôi với ánh mắt lo âu: ``Liệu chúng ta có phải về nhà mà chưa hoàn thành thủ tục không anh Kuon? Em... em không muốn thấy cảnh em ấy không được thừa nhận đâu.''

Ryuuhou cũng đứng sát lại bên tôi, giọng nhỏ vừa đủ cho mấy anh em nghe: ``Anh ơi, nếu có vấn đề gì thì mình thử hỏi bố mẹ có cách giải quyết không? Hay là\dots\ thử xem anh có quen ai làm xử lý mấy cái này không?''

Câu nói của em gái vừa dứt, thì đúng lúc chiếc đồng hồ đeo tay trên cổ tay tôi bỗng chốc chớp sáng liên hồi. Game - dù ban đầu còn đang trong chế độ nghỉ ngơ, lững thững xem phim hành động trong không gian ảo của mình - vẫn kịp thời bật dậy thực hiện một thao tác ``ảo thuật công nghệ'' khiến tất cả chúng ta bất ngờ. 

Màn hình máy in của trụ sở UBND phường, vốn đang im lặng từ nãy giờ, đột nhiên tự động khởi động, giấy liên tục chạy ra với một bản giấy giám định sinh trắc học đầy đủ dấu ấn y tế, in rõ ràng mọi thông tin được kết xuất trực tiếp từ dữ liệu quét của ông ta tại vũ trụ Aogami. Tất cả chỉ số sinh học trên giấy đều khẳng định Naomi đã tròn 8 tuổi, ngày tháng năm sinh được ấn định chính xác là 29/05/2013, hoàn toàn trùng khớp với thời điểm em được tái sinh vào năm 2021 ở thế giới em ấy từng sống.

Không chỉ dừng lại ở đó, chỉ trong vài giây ngắn ngủi, Game thực hiện thêm cú ``hack nhẹ'' vào toàn bộ hệ thống dữ liệu hộ tịch quốc gia, tự động chèn mã định danh hợp lệ và đầy đủ thông tin cho cả ba cô em gái vào cơ sở dữ liệu chung. Tất cả những rào cản pháp lý mà vừa nãy còn khiến vị cán bộ phải vò đầu bứt tai, đều bị ông ta dọn sạch chỉ trong vòng ba nốt nhạc.

Tôi đứng đó nhìn cảnh tượng diễn ra trước mắt, không khỏi thở phào nhẹ nhõm: ``Cũng may là có ông ta\dots\ nếu không thì tối nay mình phải về nhà mà lo lắng mất.'' Kaigou cũng thở ra một hơi, gật đầu đồng tình với tôi, còn Suzuno và Naomi thì vỗ tay reo lên nhỏ: ``Tuyệt quá! Vậy là bọn em có thể ở cùng mọi người rồi!'' Ryuuhou cũng mừng rỡ nhảy cẫng lên: ``Game đúng là cứu tinh của chúng ta hôm nay!''

Vị cán bộ phường cầm cả xấp hồ sơ trên tay, gồm kết quả xét nghiệm ADN, hai tờ giấy khai sinh của Kaigou và Suzuno từ thế giới song song, cộng thêm bản phân tích sinh trắc học ảo diệu vừa in ra của Naomi, đầu óc quay mòng mòng không kịp hiểu chuyện gì vừa xảy ra. Sau vài phút ngẩn ngơ nhìn đống giấy tờ, cuối cùng chú thở dài, cầm con dấu và gõ *CỘP* một cái thật to lên tờ khai nhập hộ khẩu, miệng không ngừng khen: ``Tuyệt vời\dots\ Gia đình anh chị đi thất lạc con cái chừng ấy năm mà giờ tìm lại đủ cả bộ thế này thì đúng là phép màu của khoa học hiện đại. Tôi làm việc ở đây mười mấy năm rồi, chưa bao giờ thấy một trường hợp kỳ diệu như thế này!''

Cả nhà Hikawa chỉ biết nhìn nhau cười trừ. Vậy là bằng sự kết hợp giữa di vật tình thân và bàn tay mờ ám của công nghệ thứ nguyên, Hikawa Kaigou, Hikawa Suzuno và Hikawa Naomi đã chính thức trở thành những công dân hợp pháp, có tên trong cùng một sổ hộ khẩu với tôi.

%==========================================TRUYỆN PHỤ=========================================
\omk

Ngày ngày trôi qua, ngôi nhà Hikawa bỗng chốc chuyển từ sự yên ấm thường nhật thành một bầu trời náo nhiệt chưa từng có, tiếng cười nói vang vọng từ phòng bếp đến sân sau, mọi ngóc ngách đều đong đầy hơi ấm của một gia đình thực sự. Tôi - Kuon, người anh cả mới được ``thăng chức'' chưa đầy một tuần - vẫn đôi khi đứng ngẩn người giữa căn phòng đầy ắp tiếng nói, không dám tin số phận lại xếp đặt mọi thứ tuyệt vời đến thế.

Kaigou, em song sinh của tôi, dù vẫn giữ được vẻ ngoài lạnh lùng, khó gần với những người ngoài, như lần nào gặp đám bạn đồng nghiệp ở công ty cũng chỉ chào hỏi qua loa, không bao giờ thể hiện cảm xúc gì. Nhưng mỗi khi quay về nhà, cất tiếng gọi ``Aniki'' hay ``anh Kuon'' với tôi, giọng cô ấy lại dịu đi hẳn, như gió thu thổi qua mặt hồ yên ả - vẫn còn cái mát lạnh đặc trưng, nhưng lại trong vắt, mang theo sự tin tưởng không thể che giấu. 

Tôi còn nhớ có lần đang ngồi làm việc, lỡ tay xóa nhầm cả một tập tin mã nguồn quan trọng của cô ấy, tôi đã đứng sững sờ chờ đợi cơn giận dữ bùng nổ, thế mà Kaigou chỉ thở dài, gõ nhẹ vào đầu tôi: ``Lần sau cẩn thận hơn đi, anh này lúc nào cũng lơ đễnh thế. May là em đã lưu bản sao đám mây, không thì hôm nay cả hai phải thức trắng đêm để viết lại đấy.''

Thậm chí có đêm tôi thức dậy giữa khuya đi lấy nước, vẫn bắt gặp bóng dáng cao ráo của cô ấy đang đứng ở bàn làm việc của tôi, nhẹ nhàng xếp gọn những đống báo cáo đồ án tôi để bừa bộn, để gọn máy tính rồi mới lặng lẽ quay về phòng của mình. Tôi đứng ở cửa cầu thang, nhìn theo bóng dáng em, lòng ấm áp đến muốn rơi nước mắt.

Suzuno thì hoàn toàn khác, em hệt như một con mèo nhỏ ngoan ngoãn, luôn lẻn đi sau mọi người trong nhà, chỉ muốn được ở gần mẹ hoặc các anh chị. Điều tuyệt vời nhất là bây giờ em đã chính thức học chung trường cấp hai với Ryuuhou, mỗi sáng hai chị em cùng nhau đạp xe đi học, chiều về cũng cùng nhau mang theo túi bánh mì ăn lót dạ, nói chuyện ríu rít suốt đường. 

Sáng nào cũng vậy, Suzuno là người đầu tiên thò đầu vào bếp, đôi mắt hai màu lấp lánh như ngọc nhìn mẹ đang nấu ăn, hỏi nhỏ: ``Mẹ ơi, hôm nay có món canh miso con thích không ạ? Hôm qua học về công thức nấu canh ở nhà trường, con muốn thử làm cùng mẹ!'' Cái tiếng ``Mẹ'' ấy thốt ra ngọt ngào, đong đầy bao nhiêu năm mong mỏi ấp ủ, đến nỗi mẹ tôi mỗi lần nghe đều phải quay đi lau khóe mắt, rồi xoa đầu em mà đáp: ``Có chứ, mẹ đã chuẩn bị nguyên liệu từ tối qua rồi, chúng ta cùng làm nhé.''

Còn Naomi, cô em út của chúng ta, bây giờ đã chính thức mang tên Hikawa Naomi, được nhập học vào ngôi trường tiểu học mà tôi và Ryuuhou từng theo học. Ngày đầu tiên đi học, em mặc bộ đồng phục mới tinh, đeo cặp sách nhỏ màu hồng, chạy vòng quanh nhà để khoe với mọi người, tíu tít theo sau chân bố hoặc các anh chị như một thiên thần nhỏ. Tiếng cười trong trẻo, hồn nhiên của cô bé 8 tuổi vang vọng khắp mọi nơi, xua tan đi mọi muộn phiền mệt mỏi sau những giờ làm việc, học tập căng thẳng. Mỗi chiều tôi (hoặc Ryuuhou nếu tôi quá bận) đón em từ trường về, em vừa chạy ra vừa ôm chầm lấy chân tôi, kể đủ thứ chuyện ở lớp: hôm nay cô giáo khen em vẽ tranh đẹp, bạn cùng lớp tặng em kẹo dâu, tất cả những điều nhỏ nhặt nhất em cũng đều muốn chia sẻ với anh cả.

Buổi tối nào cũng vậy, Coco và Aloe thường xuyên lên nhà chơi, lại cùng Ryuuhou, Suzuno và Naomi ngồi vây quanh bàn ăn lớn trong bếp, cùng nhau làm bánh ngọt, tiếng nói cười không ngớt. Nàng Rồng, người từng là người đầu tiên trêu chọc chúng ta hồi mới gặp, nói rằng ``Nhìn lại thì đúng là giống anh em song sinh mất rồi, thế mà lại không phải ruột thịt'', bây giờ mỗi lần đến nhà lại vừa nhào bột vừa cười lớn: ``Thế mà có ngờ đâu lời đùa của em lại thành sự thật thật! Hôm nào thấy Kaigou-paisen gọi Kuon-paisen là Aniki, em còn phải ngạc nhiên lâu, tưởng mình nghe nhầm. Thật sự là số phận sắp đặt hết rồi!''

Aloe bên cạnh cũng gật đầu cười, vừa trang trí lớp kem lên chiếc bánh vừa nói: ``Đúng vậy, hồi đầu gặp ba em Kaigou, Suzuno, Naomi, ai cũng nghĩ chỉ là bạn bè đồng hành thôi, ai ngờ bây giờ lại thành một gia đình hoàn chỉnh. Nhà Hikawa bây giờ mới thực sự đầy đủ tiếng cười.''

% Còn Kaigou thì luôn ngồi cạnh tôi trên ghế dài ở phòng khách, im lặng phụ tôi rà soát lại các dòng lệnh lập trình, có chỗ nào sai sót em chỉ khẽ chỉ vào màn hình, nói nhỏ: ``Ở dòng này, anh viết nhầm tham số rồi, sửa lại thành 0x45 mới đúng.'' Tôi quay sang nhìn em, má em hơi ửng đỏ khi bắt gặp ánh mắt của tôi, vội quay lại nhìn màn hình như đang tập trung cao độ, nhưng tôi thấy được nụ cười nhỏ vơi trên môi em.

Nhìn lại bốn cô em gái của mình, họ đều mang theo một cái chất riêng của mình: Kaigou người bản lĩnh, cứng cỏi nhưng luôn biết quan tâm mọi người; Ryuuhou cô em tomboy hoạt bát, lúc nào cũng tràn đầy năng lượng; Suzuno cô bé nhút nhát nhưng lịch sự và thông minh, yêu thích hóa học như một phù thủy nhỏ; và Naomi cô em út thiên thần, luôn mang lại niềm vui cho tất cả. Chỉ trong vòng một tuần ngắn ngủi, tôi từ một người anh chỉ có mỗi em gái Ryuuhou, bỗng chốc trở thành anh cả của bốn người, có thêm những người em gái ruột thịt để chăm sóc, để bảo vệ.

Tôi lau nhẹ giọt nước mắt vui mừng ở khóe mắt, nhìn mọi người đang vui vẻ làm bánh trong bếp, hiểu ra một điều: gia đình hóa ra không chỉ là nơi chúng ta sinh ra, mà còn là nơi những dòng máu và linh hồn tìm thấy nhau giữa vô tận các chiều không gian, vượt qua mọi bao la để về chung một mái nhà.

\end{document}